% Worked solutions for Listas/lista-topic4-moral-hazard-overhang.md (15 exercises).
% Compile: pdflatex twice; check: pdffonts (Type 1 only, uni = yes).
% Every number is asserted in lista-topic4-moral-hazard-overhang-solutions_check.py.
\documentclass[11pt,a4paper]{article}
\usepackage[utf8]{inputenc}
\usepackage[T1]{fontenc}
\usepackage{lmodern}   % Type 1 vector fonts (replaces PK bitmaps)
\usepackage{cmap}      % ToUnicode CMap -> copyable text and accents
\usepackage[brazilian,english]{babel}
\usepackage{amsmath,amssymb,amsfonts}
\usepackage[a4paper,margin=2.2cm]{geometry}
\usepackage{booktabs}
\usepackage{enumitem}
\usepackage{fancyhdr}
\usepackage{xcolor}
\usepackage{microtype}

\pagestyle{fancy}\fancyhf{}
\lhead{\footnotesize Finan\c{c}as II --- Corporativas (MPE Insper)}
\rhead{\footnotesize Topic 4 list --- Worked solutions}
\cfoot{\footnotesize\thepage}
\renewcommand{\headrulewidth}{0.3pt}
\setlength{\headheight}{14pt}
\setlist{topsep=2pt,itemsep=1pt}

\newcommand{\exsec}[1]{\section*{\color{blue!55!black}#1}\addcontentsline{toc}{section}{#1}}
\newcommand{\notation}{\par\noindent\textbf{Notation.}}
\newcommand{\step}[2]{\par\medskip\noindent\textbf{Step #1 --- #2.}~}
\newcommand{\intuition}{\par\medskip\noindent\textbf{Intuition.}~}
\newcommand{\bookans}{\par\medskip\noindent\textbf{Book answer.}~}
\newcommand{\refbib}{\par\smallskip\noindent\textbf{Ref:}~}
\newcommand{\assume}{\par\smallskip\noindent\textbf{Assumption.}~}
\newcommand{\pt}{\foreignlanguage{brazilian}}
\newcommand{\asked}{\par\medskip\noindent\textbf{What is asked.}~}
\newcommand{\approach}{\par\smallskip\noindent\textbf{Approach.}~}
\newcommand{\answer}[1]{\par\smallskip\noindent\fbox{\textbf{Answer #1.}}~}
\newcommand{\pitfalls}{\par\medskip\noindent\textbf{Common mistakes.}}

\usepackage[hyphens]{url}
\begin{document}
\begin{center}{\Large\bfseries Exercise list --- Topic 4: Moral hazard, credit rationing, asset substitution, debt overhang}\\[4pt]{\large Worked solutions}\end{center}

\noindent Statements are in \url{Listas/lista-topic4-moral-hazard-overhang.md} (and its PDF); each exercise is cited by number and printed page and is not restated. Short answers: \url{lista-topic4-moral-hazard-overhang-solutions-short.pdf}. Editions are the local copies: Tirole, \emph{The Theory of Corporate Finance} (2005); Amaro de Matos, \emph{Theoretical Foundations of Corporate Finance} (2001); Berk = Berk, DeMarzo \& Harford, \emph{Fundamentals of Corporate Finance}, 2nd ed. (2011); Grinblatt--Titman, \emph{Financial Markets and Corporate Strategy}, 2nd ed. (2002); Ross = Ross, Westerfield, Jordan \& Lamb, \emph{Fundamentos de Administra\c{c}\~ao Financeira}, 9th ed. (2013).

\tableofcontents

% =============================================================================
% Topic 4 list, exercises 1--8: worked solutions (fragment; label prefix t4e<n>-)
% Amounts in \$ million unless stated otherwise.
% =============================================================================

\exsec{1. DIP financing and existing creditors --- Grinblatt \& Titman 16.14, p.~593}

\asked A single conceptual question. Chapter~11 lets a bankrupt firm borrow new money that is
paid \emph{before} the old debt. On its face this hurts the old creditors, who move down the
queue. Explain the mechanism by which they can nevertheless end up better off.

\approach This is the lecture's \emph{debt overhang} problem (Myers 1977). Write every claim as
a state-by-state payoff, debt $=\min\{\text{assets},F\}$ and equity $=\max\{\text{assets}-F,0\}$,
and compare the old creditors' payoff with and without the new project. Junior money cannot be
raised when the old debt is under water, because the old creditors capture the project's value;
seniority for the new money removes that obstacle. The old creditors lose priority but gain the
project's cash flows, so the comparison is ``smaller share of a larger pie'' versus ``larger share
of a smaller pie''.

\notation
\begin{itemize}
  \item $s$: state of the world at the payoff date; $V_s\ge 0$: value of the assets in place in state $s$.
  \item $F>0$: face value of the existing debt, protected by a covenant that forbids issuing debt senior to it.
  \item $I>0$: cost of a new project that must be carried out inside the firm; $X_s\ge 0$: its payoff in state $s$; risk-neutral investors, zero riskless rate (as in the book's and the lecture's examples).
  \item $N$: face value of the new debtor-in-possession (DIP) debt, senior to $F$; $D$, $E$: market values of the existing debt and of equity.
\end{itemize}

\step{1}{Why the project is not funded outside bankruptcy} We first show that, without a
seniority exception, a good project is passed up; this is the status quo the old creditors
compare against. Without DIP financing, new money can only be junior to $F$. In state $s$ a junior claimant receives at most what is left after the
senior debt,
\begin{align}
  \text{junior payoff}_s &\le \max\{V_s+X_s-F,\,0\}. \label{t4e1-junior}
\end{align}
In the book's Lily Pharmaceuticals example (pp.~564--566), in the unfavourable state the
assets in place are worth nothing without further development ($V_s=0$), development costs
$I=100$ and yields $X_s=150$, and the senior obligation is $F=110$. Substituting into
\eqref{t4e1-junior},
\begin{align}
  \text{junior payoff} &\le \max\{(0)+(150)-(110),\,0\} \nonumber\\
                       &= \max\{40,\,0\} \nonumber\\
                       &= 40 < 100 = I, \nonumber
\end{align}
so no one lends: junior lenders would pay $100$ for a claim worth at most $40$. Yet the project
is worth doing,
\begin{align}
  \mathit{NPV} &= X_s-I \nonumber\\
               &= (150)-(100) \nonumber\\
               &= 50>0, \nonumber
\end{align}
so it is passed up, the firm defaults, and the existing lenders receive
$\min\{V_s,F\}=\min\{0,110\}=0$ in that state.

\step{2}{DIP debt makes it fundable} We now let the new money jump the queue and recompute the
old creditors' payoff. In Chapter~11 the firm may issue debt senior to $F$. A DIP
loan of $I=100$ at the zero riskless rate has face $N=100$ and is riskless because
$150\ge 100$. The existing lenders now receive
\begin{align}
  \text{existing debt payoff} &= \min\{V_s+X_s-N,\,F\} \label{t4e1-old}\\
                              &= \min\{(0)+(150)-(100),\,(110)\} \nonumber\\
                              &= \min\{50,\,110\} \nonumber\\
                              &= 50, \nonumber
\end{align}
instead of $0$: they gain $50-0=50$, the whole NPV, even though they lost their priority.

\step{3}{The general argument} The example is not special; we show the old creditors never lose
when the DIP money funds a project that pays back at least its cost. With a riskless DIP loan ($N=I\le X_s$ in every state), the
existing debt's payoff with the project is $\min\{V_s+X_s-I,\,F\}$, and without it
$\min\{V_s,\,F\}$. Since $X_s-I\ge 0$, the first argument of the minimum is larger with the
project, $V_s+X_s-I\ge V_s$, and $\min\{\cdot,F\}$ is nondecreasing in its first argument, so
\begin{align}
  \min\{V_s+X_s-I,\,F\} &\ge \min\{V_s,\,F\} \qquad\text{for every } s, \label{t4e1-dom}
\end{align}
with strict inequality in every state where the old debt was impaired ($V_s<F$). Seniority is
diluted, but the asset pool the old creditors are paid from grows by more than the new claim.

\step{4}{The lecture's numbers} Topic~4B (slides 22--31): assets in place worth $500$ (boom)
or $200$ (slump), each with probability $0.5$; $F=300$; a project costs $I=60$ and adds $100$
in both states. We value both claims without the project, then with it. Without the project,
\begin{align}
  D &= (0.5)\min\{500,300\}+(0.5)\min\{200,300\} \nonumber\\
    &= (0.5)(300)+(0.5)(200) = 150+100 = 250, \nonumber\\
  E &= (0.5)\max\{500-300,0\}+(0.5)\max\{200-300,0\} \nonumber\\
    &= (0.5)(200)+(0.5)(0) = 100. \nonumber
\end{align}
Financing the $60$ with senior DIP debt ($N=60$, riskless since the slump assets $200+100=300\ge 60$),
the assets become $500+100=600$ (boom) and $200+100=300$ (slump), and the old debt
receives, by \eqref{t4e1-old},
\begin{align}
  \text{boom: } & \min\{(600)-(60),\,(300)\} = \min\{540,300\} = 300, \nonumber\\
  \text{slump: } & \min\{(300)-(60),\,(300)\} = \min\{240,300\} = 240, \nonumber\\
  D &= (0.5)(300)+(0.5)(240) \label{t4e1-D}\\
    &= 150+120 = 270, \nonumber
\end{align}
and equity receives $\max\{600-60-300,0\}=240$ (boom) and $\max\{300-60-300,0\}=0$ (slump), so
$E=(0.5)(240)+(0.5)(0)=120$. Both claims rise by $20$ ($270-250$ and $120-100$), and
$20+20=40=100-60$ is the project's NPV.

\answer{(16.14)} The old creditors benefit because, without senior new money, a positive-NPV project
cannot be financed (debt overhang) and they are paid out of the assets in place only. With DIP
financing the project is done; the DIP lender is repaid first, but the project adds more to the
pool than the DIP claim takes out ($X_s\ge I$), so the old creditors' payoff rises in every
impaired state (\$0 to \$50 in the Lily example; \$250 to \$270 in the lecture's example).

\step{5}{The limit} \eqref{t4e1-dom} needs $X_s\ge I$: a good, safe project. The same provision
lets a distressed firm borrow to keep operating when it would be worth more liquidated; then
$X_s<I$ in some state and the old creditors lose (Result~16.11).

\intuition Debt overhang means the old creditors capture part of the value of new investment
while the new money bears its full cost, so junior money cannot be raised. Letting new money
jump the queue removes the overhang. The old creditors, who would otherwise get little or
nothing, take a smaller slice of a larger pie. A covenant that blocks this hurts them, and a
diffuse group of bondholders cannot easily agree to lend the money themselves (free-rider
problem). Chapter~11 unlocks senior financing only in distress, which keeps most of the
protection the covenant gives in normal times.
\pitfalls
\begin{itemize}
  \item Comparing the old creditors' payoff with the project under DIP financing against their payoff
        with the project under \emph{junior} financing. The relevant status quo is no project at all,
        because junior money cannot be raised.
  \item Crediting the old creditors with the project's gross payoff $X_s$ and forgetting that the DIP
        lender is repaid $N=I$ first; their gain is bounded by $X_s-I$, the project's NPV.
  \item Stating that DIP financing always helps creditors. It helps only when the funded project pays
        back at least its cost ($X_s\ge I$); financing a value-destroying continuation hurts them.
\end{itemize}
\refbib{Grinblatt \& Titman Ch.~16, \S~16.2 ``Underinvestment as a Consequence of Debt Overhang'' (Lily example), pp.~564--566; \S~16.3 ``How Chapter 11 Bankruptcy Mitigates Debt Holder--Equity Holder Incentive Problems'', Result~16.11, p.~580. Lecture Topic~4B, slides 22--31.}

% =============================================================================
\exsec{2. Asset substitution: project A or B --- Grinblatt \& Titman 16.1, p.~591}

\asked (a) The net present value of each project for the firm as a whole. (b) Which project the
shareholders, who control the firm, would pick, and the reason for their choice.

\approach With risk-neutral investors and a zero riskless rate, a claim's value is its expected
payoff. Debt and equity split each state's payoff: debt gets $\min\{X,F\}$, equity gets
$\max\{X-F,0\}$ (limited liability). Because the cash on hand exactly equals the face value, equity
is an at-the-money call option on the project payoff, so it values variance. This is the lecture's
\emph{asset substitution} (risk-shifting) problem: comparing firm value (a) with equity value (b)
shows the gap between what is efficient and what shareholders choose.

\notation
\begin{itemize}
  \item $j\in\{A,B\}$: project; $X_{j,s}$: its payoff in state $s\in\{\text{fav},\text{unf}\}$, each with probability $\pi=0.5$.
  \item $C=100$: cash on hand, all spent on the project; $F=100$: debt due next period; $r_f=0$; risk-neutral investors.
  \item $V_j=\mathbb{E}[X_j]$: firm value under $j$; $E_j$, $D_j$: values of equity and debt under $j$.
\end{itemize}

\step{1}{(a) NPVs} We value each project's payoff and subtract the \$100 it costs. With $r_f=0$
and risk neutrality the present value is the expected payoff:
\begin{align}
  V_A &= \pi X_{A,\text{fav}}+(1-\pi)X_{A,\text{unf}} \label{t4e2-VA}\\
      &= (0.5)(120)+(0.5)(60) \nonumber\\
      &= 60+30 = 90, \nonumber\\
  \mathit{NPV}_A &= V_A-C \label{t4e2-NA}\\
      &= (90)-(100) \nonumber\\
      &= -10, \nonumber\\
  V_B &= \pi X_{B,\text{fav}}+(1-\pi)X_{B,\text{unf}} \label{t4e2-VB}\\
      &= (0.5)(101)+(0.5)(101) \nonumber\\
      &= 50.5+50.5 = 101, \nonumber\\
  \mathit{NPV}_B &= V_B-C \nonumber\\
      &= (101)-(100) \nonumber\\
      &= 1. \nonumber
\end{align}

\answer{(a)} $\mathit{NPV}_A=-\$10$ million, $\mathit{NPV}_B=+\$1$ million: only B creates value.

\step{2}{(b) Payoffs to each claim, state by state} The shareholders care about their own claim,
not about firm value, so we split each project's payoff between the claims. Equity receives
$\max\{X-F,0\}$ and debt $\min\{X,F\}$:
\begin{center}
\begin{tabular}{lcccc}
\toprule
 & \multicolumn{2}{c}{Equity} & \multicolumn{2}{c}{Debt}\\
 & Favourable & Unfavourable & Favourable & Unfavourable\\
\midrule
A & $120-100=20$ & $0$ & $100$ & $60$\\
B & $101-100=1$ & $101-100=1$ & $100$ & $100$\\
\bottomrule
\end{tabular}
\end{center}
(In A's unfavourable state $\max\{60-100,0\}=0$ and $\min\{60,100\}=60$.) Taking expectations,
\begin{align}
  E_A &= (0.5)(20)+(0.5)(0) = 10+0 = 10, & D_A &= (0.5)(100)+(0.5)(60) = 50+30 = 80, \label{t4e2-A}\\
  E_B &= (0.5)(1)+(0.5)(1) = 0.5+0.5 = 1, & D_B &= (0.5)(100)+(0.5)(100) = 100. \label{t4e2-B}
\end{align}
Check: $E_A+D_A=10+80=90=V_A$ and $E_B+D_B=1+100=101=V_B$.

\step{3}{(b) The choice} We compare the shareholders' payoffs and decompose the difference into a
value effect and a transfer. Equity holders, who control the firm, want project A: $E_A=10>E_B=1$,
although A destroys value. Writing $E_j=V_j-D_j$ and using \eqref{t4e2-A}--\eqref{t4e2-B},
\begin{align}
  E_A-E_B &= (V_A-V_B)-(D_A-D_B) \label{t4e2-split}\\
          &= (90-101)-(80-100) \nonumber\\
          &= (-11)-(-20) \nonumber\\
          &= 9: \nonumber
\end{align}
switching to A lowers firm value by $11$ but moves $20$ from the creditors to the shareholders.

\answer{(b)} The shareholders want project A ($E_A=\$10$ million $>E_B=\$1$ million), because
limited liability lets them keep A's upside while the creditors bear its downside: the \$20 million
transferred from debt exceeds the \$11 million of value destroyed.

\intuition With \$100 due and \$100 in cash, the equity is an at-the-money call on the firm's
assets. The safe project leaves it almost worthless ($1$); the risky one gives it the whole
upside ($20$) while the creditors absorb the downside ($60$ instead of $100$). This is the
lecture's asset-substitution (risk-shifting) problem: limited liability makes levered equity
prefer variance, even at negative NPV.
\pitfalls
\begin{itemize}
  \item Writing equity's payoff as $X-F$ in every state. In A's unfavourable state that gives $-40$;
        limited liability caps the loss at $0$, and that cap is exactly what makes A attractive.
  \item Answering (b) with the NPVs of (a). Firm value and equity value rank the projects in opposite
        orders; (b) asks about the shareholders' claim.
  \item Forgetting that the \$100 of cash is spent on the project: the NPV subtracts it, but the
        payoff split in (b) uses only the project payoff and the \$100 face value.
\end{itemize}
\refbib{Grinblatt \& Titman Ch.~16, \S~16.2 ``The Asset Substitution Problem'', pp.~569--572 (Results~16.7--16.8). Lecture Topic~4B, slides 2--8.}

% =============================================================================
\exsec{3. Debt overhang on vacant land --- Berk 16.25, pp.~494--495}

\asked (a) Today's values of equity and debt if the land stays vacant. (b) The NPV of developing
the land, for the firm as a whole. (c) Today's values of equity and debt if the shareholders put in
\$20 million and the land is developed. (d) Whether the shareholders gain from putting in that
money, given (c).

\approach All cash flows are riskless, so every claim is its payoff in one year discounted at
$r_f=10\%$; debt gets $\min\{\text{land value},15\}$ and equity the rest. This is the lecture's
\emph{debt overhang} (under-investment) problem: the debt is under water, so part of the value a
new project creates goes to the creditors, who pay none of its cost. Comparing the NPV (b) with the
shareholders' gain (d) shows the wedge.

\notation
\begin{itemize}
  \item $F=15$: face value of the debt, due in one year; $r_f=0.10$; all cash flows riskless; no taxes.
  \item $L_0=10$: value of the vacant land in one year; $L_1=35$: value of the developed land in one year; $K=20$: development cost paid today.
  \item $E$, $D$: values today of equity and debt.
\end{itemize}

\step{1}{(a) Land left vacant} We split the land's value in one year between the claims and
discount. In one year the firm is worth $10<15$: the creditors take $10$ and the shareholders get
$0$. Discounting at $r_f$,
\begin{align}
  D &= \frac{\min\{L_0,F\}}{1+r_f} \label{t4e3-Da}\\
    &= \frac{\min\{10,15\}}{1+0.10} \nonumber\\
    &= \frac{(10)}{(1.10)} \nonumber\\
    &= 9.0909, \nonumber\\
  E &= \frac{\max\{L_0-F,0\}}{1+r_f} = \frac{\max\{10-15,0\}}{1.10} = \frac{(0)}{(1.10)} = 0. \nonumber
\end{align}

\answer{(a)} Equity is worth \$0; debt is worth \$9.09 million.

\step{2}{(b) NPV of developing} The NPV counts only the cash flows that change with the decision.
Developing replaces land worth $10$ by land worth $35$, so the incremental cash flow in one year is
$35-10=25$, against $20$ paid today:
\begin{align}
  \mathit{NPV} &= -K+\frac{L_1-L_0}{1+r_f} \label{t4e3-NPV}\\
      &= -(20)+\frac{(35)-(10)}{(1.10)} \nonumber\\
      &= -20+\frac{25}{1.10} \nonumber\\
      &= -20+22.7273 \nonumber\\
      &= 2.7273. \nonumber
\end{align}
(Ignoring the vacant-land value forgone would give $-20+35/1.10=11.8182$, the NPV measured
against a firm with no land at all; the decision-relevant figure is \eqref{t4e3-NPV}.)

\answer{(b)} $\mathit{NPV}=\$2.73$ million $>0$: developing creates value.

\step{3}{(c) Values if developed with \$20 of new equity} We repeat Step 1 with the developed land.
The \$20 is spent on development, so it is not an asset in one year. In one year the firm is worth
$35\ge 15$: the debt is paid in full, $\min\{35,15\}=15$, and equity keeps $\max\{35-15,0\}=20$. Hence
\begin{align}
  D &= \frac{(15)}{(1.10)} \label{t4e3-Dc}\\
    &= 13.6364, \nonumber\\
  E &= \frac{(35)-(15)}{(1.10)} \label{t4e3-Ec}\\
    &= \frac{20}{1.10} \nonumber\\
    &= 18.1818. \nonumber
\end{align}

\answer{(c)} Equity is worth \$18.18 million and debt \$13.64 million today.

\step{4}{(d) Will the shareholders put up \$20?} The shareholders compare what they hold after
contributing with what they hold if they do nothing. Their gain is the equity value
\eqref{t4e3-Ec}, minus the equity value without development (Step 1), minus the contribution:
\begin{align}
  \Delta E &= E_{\text{dev}}-E_{\text{vacant}}-K \nonumber\\
           &= (18.1818)-(0)-(20) \nonumber\\
           &= -1.8182 < 0, \nonumber\\
  \Delta D &= D_{\text{dev}}-D_{\text{vacant}} \nonumber\\
           &= (13.6364)-(9.0909) \nonumber\\
           &= 4.5455. \nonumber
\end{align}
No: they would lose \$1.82 million. The sum $-1.8182+4.5455=2.7273$ is the NPV
\eqref{t4e3-NPV}; the creditors capture more than all of it.

\answer{(d)} No. The shareholders would pay \$20 million for an equity stake worth \$18.18 million
(a loss of \$1.82 million), while the debt gains \$4.55 million; the positive-NPV project is passed up.

\intuition The debt is under water ($10<15$), so the first \$5 million of any value the
project creates goes to the creditors, who pay none of its cost. The shareholders bear the
\$20 million and receive only the value above the \$15 million face. This is Myers' debt
overhang: leverage acts as a tax on new investment, and a positive-NPV project is passed up.
Writing the debt down (renegotiation, as in the lecture) or senior new financing would fix it.
\pitfalls
\begin{itemize}
  \item Computing the NPV as $-20+35/1.10=11.82$: the vacant land worth \$10 million in one year is
        given up, so it is part of the project's cost.
  \item Answering (d) with $E_{\text{dev}}=18.18>0$, ``so yes''. The shareholders must compare the
        stake with the \$20 million they put in.
  \item Valuing the debt in (a) at the face value $15/1.10$; the creditors can get at most what the
        land is worth, $\min\{10,15\}$.
\end{itemize}
\refbib{Berk Ch.~16, \S~16.6 ``Agency Costs and Information'', ``Equity-Debt Holder Conflicts'' (under-investment), pp.~483--484. Lecture Topic~4B, slides 19--30.}

% =============================================================================
\exsec{4. Short-term debt and incentive problems --- Grinblatt \& Titman 16.11, p.~593}

\asked A single conceptual question: why the conflicts between shareholders and creditors
(asset substitution, debt overhang, under-maintenance) do less damage when the firm finances itself
with debt that matures soon and must be rolled over, rather than with long-term debt.

\approach The conflicts exist because equity chooses the firm's strategy \emph{after} the terms of
the debt are fixed, so it can change the value of the creditors' claim (the lecture's asset
substitution and debt overhang). We write equity as firm value minus debt value and ask what
happens to the debt term when the debt is repriced after the choice is made: if new lenders price
the debt knowing the strategy, the debt term no longer depends on the strategy and equity
internalizes the full value consequences.

\notation
\begin{itemize}
  \item $j$: investment strategy (project choice, risk level, maintenance), chosen by the equity holders after the debt is in place; $X_j\ge 0$: resulting terminal firm value, $V_j=\mathbb{E}[X_j]$ (risk neutrality, $r_f=0$).
  \item $F$: face value of the debt; $L$: amount the firm owes when its short-term debt is rolled over; $D_j$, $E_j$: values of debt and equity under $j$.
\end{itemize}

\step{1}{Where the conflict comes from} We locate the conflict in the debt term of equity's value.
With long-term debt, $F$ is fixed before $j$ is chosen, and since $\max\{X-F,0\}=X-\min\{X,F\}$ in
every state,
\begin{align}
  E_j &= \mathbb{E}[\max\{X_j-F,0\}] \label{t4e4-E}\\
      &= V_j-D_j, \qquad D_j=\mathbb{E}[\min\{X_j,F\}]. \nonumber
\end{align}
Equity can gain by picking a $j$ that lowers $D_j$ (more risk, less investment, payouts) even
when $V_j$ falls: comparing two strategies, $E_j-E_k=(V_j-V_k)-(D_j-D_k)$, and the second term
is a transfer from creditors whenever $D_j<D_k$.

\step{2}{Debt that is repriced when the choice is made} We now let the debt be repriced after the
choice and show the transfer term disappears. If the debt matures before or at the
decision and must be rolled over, the new lenders see $j$ and set the new face $F_j$ so that
their claim is worth what they lend, $\mathbb{E}[\min\{X_j,F_j\}]=L$. Substituting into
\eqref{t4e4-E},
\begin{align}
  E_j &= V_j-(L), \label{t4e4-Est}
\end{align}
and $L$ does not depend on $j$: equity bears the full consequences of its choice and picks
the value-maximizing strategy. Short-term debt is worth close to its face whatever the firm
does, so there is little value to expropriate.

\step{3}{Illustration (Grinblatt \& Titman 16.1)} Take the firm of G\&T Exercise~16.1: \$100 of
debt, projects A (pays $120$ or $60$) and B (pays $101$ in both states), each state with
probability $0.5$, $r_f=0$. With long-term debt, equity is worth $0.5(20)+0.5(0)=10$ under A and
$0.5(1)+0.5(1)=1$ under B, so it picks the value-destroying A. Suppose instead the \$100 had to be
refinanced after the project is chosen. Firm values are
$V_A=0.5(120)+0.5(60)=90$ and $V_B=0.5(101)+0.5(101)=101$; by \eqref{t4e4-Est} with $L=100$,
\begin{align}
  E_A &= (90)-(100) = -10, & E_B &= (101)-(100) = 1. \nonumber
\end{align}
Under A the firm cannot raise \$100 (no claim on payoffs worth $90$ is worth $100$), so the
equity holders choose B, the positive-NPV project, instead of A.

\step{4}{Debt overhang} Myers (1977): if the existing debt matures before the firm must raise
money for a new project, the debt is repriced together with the project decision. The
project's benefit to creditors (a safer firm) then shows up as a lower borrowing rate that
accrues to equity, and positive-NPV projects are taken. In Unitron's case, the riskless
12\% project would lower the firm's 13\% borrowing rate and become attractive.

\step{5}{Asset substitution} The lending rate is renegotiated after the risk choice is made and
seen, so a risky choice raises the firm's own cost of funds instead of transferring wealth
from existing lenders. The threat of not being rolled over also disciplines the choice (lecture
Topic~4B, slide~16).

\step{6}{The cost} Rolling over exposes the firm to refinancing and interest-rate risk: an
unexpected rise in rates can push a highly levered firm into distress (hence the advice to
pair short-term debt with interest-rate hedges).

\answer{(16.11)} Short-term debt is repriced (rolled over) after the shareholders' investment and
risk choices are made and observed, so its terms reflect those choices: equity bears the full value
consequences, $E_j=V_j-L$, and cannot shift wealth from creditors through risk-taking or
under-investment. Short-term debt is also close to riskless, so there is little value to
expropriate. The price is rollover and interest-rate risk.

\intuition The conflicts exist because equity can change the value of a claim whose terms are
already fixed. The shorter the maturity, the sooner those terms are reset at prices that
reflect what equity has done, and the less wealth equity can take from creditors. The
incentive problem shrinks, and the firm takes on rollover risk in exchange.
\pitfalls
\begin{itemize}
  \item Saying short-term debt helps ``because it is safer''. The point is timing: the debt is priced
        after the decision, so the decision's effect on the creditors comes back to equity through the
        new rate.
  \item Confusing maturity with seniority. Seniority decides who is paid first in default; maturity
        decides when the terms are reset.
  \item Omitting the cost. Short maturities expose the firm to refinancing and interest-rate risk, so
        they are not a free cure.
\end{itemize}
\refbib{Grinblatt \& Titman Ch.~16, \S~16.4 ``The Use of Short-Term versus Long-Term Debt'', pp.~584--585 (Result~16.12). Lecture Topic~4B, slide~16.}

% =============================================================================
\exsec{5. Random financing --- Tirole 3.1, p.~144}

\asked (i) Write the condition under which investors are willing to sign a contract in which the
borrower contributes $\hat A$ of her cash, the project then goes ahead only with probability $x$,
and she gets $R_b$ on success. (ii) Show that a borrower with too little cash for certain financing
($A<\bar A$) should put \emph{all} her cash in, and describe how the financing probability depends
on $A$.

\approach This is the Holmstr\"om--Tirole fixed-investment model of the lecture (Topic~4A). The two
constraints are the borrower's incentive constraint (IC$_b$), which caps the \emph{pledgeable income}
at $\mathcal P=p_H(R-B/\Delta p)$, and the investors' breakeven constraint (IR$_l$). Below $\bar A$
the project cannot be financed for sure; a lottery over financing scales the investors' outlay and
the pledgeable income by $x$ but not the borrower's contribution, which convexifies the
indivisible project. We write (IR$_l$) for the lottery, show (IC$_b$) must hold, and maximise the
borrower's utility over $(x,\hat A)$.

\notation
\begin{itemize}
  \item $I>0$: fixed investment; $A\in(0,\bar A)$: borrower's cash; $\hat A\in[0,A]$: the part she contributes; $x\in[0,1]$: probability that the project is financed.
  \item $R>0$: income on success (0 on failure); $R_b\ge 0$: borrower's reward on success; $R_l=R-R_b$: investors' claim on success.
  \item $p_H>p_L$: success probabilities when she behaves / misbehaves; $\Delta p=p_H-p_L$; $B>0$: private benefit from misbehaving.
  \item $\mathcal P=p_H(R-B/\Delta p)$: pledgeable income; $\bar A=I-\mathcal P$: minimum cash for financing with probability~1 (the statement's definition); $U_b$: borrower's net utility. Rate of interest 0, everybody risk neutral.
\end{itemize}

\assume The borrower hands over $\hat A$ when the contract is signed, whether or not the
lottery then selects the project (this is the reading in Tirole's answer). If $\hat A$ were
paid only when the project goes ahead, the expected outlay would be $x(I-\hat A)$ and the
breakeven condition $x(I-\hat A)\le x\mathcal P$ would require $\hat A\ge\bar A>A$: randomizing
would achieve nothing.

\step{1}{(i) Investors' breakeven condition} Investors receive $\hat A$ for sure, pay $I$ with
probability $x$, and, if the borrower behaves, collect $R_l$ with probability $p_H$ when the
project runs:
\begin{align}
  \underbrace{\hat A}_{\text{received}}+\underbrace{x\,p_H(R-R_b)}_{\text{expected repayment}}
     &\ge \underbrace{x I}_{\text{expected outlay}} \nonumber\\
  x\,p_H(R-R_b) &\ge x I-\hat A. \label{t4e5-IR}
\end{align}

\answer{(i)} $x\,p_H(R-R_b)\ge xI-\hat A$ (with $R_b\ge B/\Delta p$ so that she behaves).

\step{2}{The contract must induce good behaviour} We check that a contract under which she shirks
is never worth signing, so (IC$_b$) can be imposed. Suppose $(\Delta p)R_b<B$, so she shirks.
Investors break even only if $x p_L(R-R_b)\ge xI-\hat A$, i.e.\ $-\hat A\le x p_L(R-R_b)-xI$.
Her net utility is then
\begin{align}
  x(p_L R_b+B)-\hat A &\le x(p_L R_b+B)+\big(x p_L(R-R_b)-xI\big) \nonumber\\
                      &= x p_L R_b+xB+x p_L R-x p_L R_b-xI \nonumber\\
                      &= x(p_L R+B-I) < 0 \label{t4e5-shirk}
\end{align}
by (3.2), worse than not signing. So the contract satisfies (IC$_b$), $(\Delta p)R_b\ge B$, i.e.\
$R_b\ge B/\Delta p$.

\step{3}{Breakeven in terms of $\bar A$} We combine the two constraints into one feasibility
condition on $(x,\hat A)$. (IC$_b$) gives $R-R_b\le R-B/\Delta p$, so the
left side of \eqref{t4e5-IR} is at most $x\,p_H(R-B/\Delta p)=x\mathcal P$. A contract exists for $(x,\hat A)$ only if
\begin{align}
  x I-\hat A &\le x\,p_H\left(R-\frac{B}{\Delta p}\right) \nonumber\\
  x I-\hat A &\le x(\mathcal P) \nonumber\\
  x(I-\mathcal P) &\le \hat A \nonumber\\
  x(\bar A) &\le \hat A \nonumber\\
  x &\le \frac{\hat A}{\bar A}, \label{t4e5-x}
\end{align}
substituting $I-\mathcal P=\bar A$ from the statement's definition (third to fourth line) and
dividing by $\bar A>0$ (last line). Conversely, if \eqref{t4e5-x} holds, $R_b=B/\Delta p$ satisfies
both constraints, so \eqref{t4e5-x} is necessary and sufficient.

\step{4}{(ii) The borrower's objective} We write her utility and use the binding breakeven
constraint to express it in terms of $x$ only. She keeps $A-\hat A$ and, with probability $x$, earns
$R_b$ with probability $p_H$; net of the $A$ she would consume anyway,
\begin{align}
  U_b &= x\,p_H R_b+(A-\hat A)-A \nonumber\\
      &= x\,p_H R_b-\hat A. \label{t4e5-Ub}
\end{align}
Competitive investors make \eqref{t4e5-IR} bind; expanding and solving for $x p_H R_b$,
\begin{align}
  x p_H R-x p_H R_b &= xI-\hat A \nonumber\\
  x p_H R_b &= x p_H R-xI+\hat A. \nonumber
\end{align}
Substituting into \eqref{t4e5-Ub},
\begin{align}
  U_b &= (x p_H R-xI+\hat A)-\hat A \nonumber\\
      &= x\,(p_H R-I). \label{t4e5-Ubx}
\end{align}

\step{5}{(ii) Optimal $\hat A$ and $x$} We maximise \eqref{t4e5-Ubx} over the feasible set
\eqref{t4e5-x}. Differentiating,
\begin{align}
  \frac{\partial U_b}{\partial x} = p_H R-I>0, \qquad \frac{\partial U_b}{\partial \hat A}=0, \nonumber
\end{align}
so $U_b$ rises with $x$ and $\hat A$ matters only through feasibility. By \eqref{t4e5-x}, the
largest feasible $x$ is $\hat A/\bar A$ (below 1 because $\hat A\le A<\bar A$), which increases with
$\hat A$: substituting it, $U_b=(\hat A/\bar A)(p_HR-I)$, increasing in $\hat A$ on $[0,A]$. So
she contributes everything:
\begin{align}
  \hat A^* &= A, \nonumber\\
  x^* &= \frac{A}{\bar A}, \label{t4e5-xstar}\\
  U_b^* &= \frac{A}{\bar A}\,(p_H R-I). \label{t4e5-Ubstar}
\end{align}
At $x^*$ the binding \eqref{t4e5-IR}, divided by $x^*=A/\bar A$, gives
$p_H(R-R_b)=I-A/x^*=I-\bar A=\mathcal P$, so $R-R_b=R-B/\Delta p$ and $R_b=B/\Delta p$: (IC$_b$) binds.

\step{6}{(ii) How $x^*$ varies with $A$} $x^*=A/\bar A$ rises linearly from 0 at $A=0$ to 1 at
$A=\bar A$, and stays at 1 above $\bar A$. Utility \eqref{t4e5-Ubstar} therefore rises smoothly
to the full NPV instead of jumping at $\bar A$ as in (3.6).

\emph{Numerical illustration.} $I=100$, $R=150$, $p_H=0.8$, $p_L=0.4$, $B=30$, $A=10$: (3.2)
holds ($0.4(150)-100+30=-10<0$); $B/\Delta p=75$; $\mathcal P=0.8(150-75)=60$; $\bar A=100-60=40$;
NPV $=0.8(150)-100=20$. Then $x^*=10/40=0.25$ and $U_b^*=0.25(20)=5$; investors pay
$0.25(100)-10=15$ in expectation and receive $0.25(60)=15$.

\answer{(ii)} $\hat A^*=A$; the financing probability is $x^*=A/\bar A$, rising linearly from $0$
at $A=0$ to $1$ at $A=\bar A$ (and $x^*=1$ for $A\ge\bar A$); her utility is $x^*(p_HR-I)$.

\intuition The obstacle is not the size of the NPV but the pledgeable income per unit of
investment. A poor borrower can still buy a \emph{fraction} of the deal: each unit of her
cash covers the investors' shortfall on $1/\bar A$ units of probability of financing. Since
she captures the whole NPV whenever the project runs, she wants as high a probability as her
cash can buy and puts all of it in. Lotteries convexify the indivisible project.
\pitfalls
\begin{itemize}
  \item Writing the breakeven as $x p_H(R-R_b)\ge x(I-\hat A)$, i.e.\ taking her contribution only when
        the project goes ahead. Then $x$ cancels and randomisation achieves nothing.
  \item Using $R$ instead of $R-B/\Delta p$ as the investors' maximal claim: without (IC$_b$) she
        shirks, and by (3.2) no one breaks even.
  \item Claiming she keeps some cash ``for safety''. Cash kept aside earns nothing extra, while each unit
        put in raises $x$ by $1/\bar A$ and her utility by $(p_HR-I)/\bar A>0$.
\end{itemize}
\bookans Matches the book's answer (p.~571): breakeven $xI-A\le x p_H(R-R_b)$, equivalently
$x\bar A\le A$; $U_b=x(p_H R-I)$, so $x^*=A/\bar A$, rising from 0 to 1 as $A$ goes from 0
to $\bar A$.
\refbib{Tirole Ch.~3, \S~3.2.1--3.2.2 ``The Fixed-Investment Model'', ``The Lenders' Credit Analysis'', pp.~115--117 (footnote~9, p.~116). Lecture Topic~4A, slides 26--34.}

% =============================================================================
\exsec{6. Zymase: leverage and strategy choice --- Berk 16.26, p.~495}

\asked (a) Which strategy gives the largest expected payoff to the firm as a whole. (b) and (c)
Which strategy gives the shareholders the largest expected payoff when \$40 million, respectively
\$110 million, of debt is due at the payoff date. (d) How much expected firm value is lost, in each
case, when management follows the shareholders' preference instead of the value-maximising strategy
(the expected agency cost of debt).

\approach Split each strategy's payoff state by state into debt, $\min\{X,F\}$, and equity,
$\max\{X-F,0\}$, and take expectations. Equity is a call option on the firm, so as $F$ rises it
favours strategies whose payoff lies above $F$ with some probability, even if their mean is lower:
the lecture's \emph{asset substitution} (excessive risk-taking). The agency cost is the drop in
expected \emph{firm} payoff caused by that choice, which is the loss the shareholders bear ex ante
when creditors price the debt.

\notation
\begin{itemize}
  \item $j\in\{A,B,C\}$: research strategy; $X_j$: its after-tax payoff (\$ million), with the probabilities in the table.
  \item $F\in\{40,110\}$: face value of the debt, due when the payoff arrives; equity gets $\max\{X_j-F,0\}$, debt $\min\{X_j,F\}$.
  \item Agency cost: expected firm payoff of the value-maximizing strategy minus that of the strategy equity chooses.
\end{itemize}

\assume Risk is diversifiable, so every claim is discounted at the riskless rate, which is not
given. All comparisons are in expected payoffs at the payoff date; present values are these
divided by the same factor $1+r_f$, so the rankings do not change.

\step{1}{(a) Expected firm payoffs} We weight each payoff by its probability:
\begin{align}
  \mathbb{E}[X_A] &= (1.0)(75) = 75, \label{t4e6-XA}\\
  \mathbb{E}[X_B] &= (0.5)(140)+(0.5)(0) = 70+0 = 70, \label{t4e6-XB}\\
  \mathbb{E}[X_C] &= (0.1)(300)+(0.9)(40) = 30+36 = 66. \label{t4e6-XC}
\end{align}
Strategy A has the highest expected payoff (and the lowest risk).

\answer{(a)} Strategy A, with an expected payoff of \$75 million (B: \$70 million; C: \$66 million).

\step{2}{(b) Debt of \$40 million} We split each state's payoff with $F=40$: equity
$\max\{X-40,0\}$ (e.g.\ $140-40=100$; $\max\{40-40,0\}=0$), debt $\min\{X,40\}$:
\begin{center}
\begin{tabular}{lccc}
\toprule
Strategy, state (prob.) & Firm & Equity $\max\{X-40,0\}$ & Debt $\min\{X,40\}$\\
\midrule
A (1.0) & 75 & 35 & 40\\
B, high (0.5) & 140 & 100 & 40\\
B, low (0.5) & 0 & 0 & 0\\
C, high (0.1) & 300 & 260 & 40\\
C, low (0.9) & 40 & 0 & 40\\
\bottomrule
\end{tabular}
\end{center}
\begin{align}
  \mathbb{E}[E_A] &= (1.0)(35) = 35, \nonumber\\
  \mathbb{E}[E_B] &= (0.5)(100)+(0.5)(0) = 50+0 = 50, \label{t4e6-EB40}\\
  \mathbb{E}[E_C] &= (0.1)(260)+(0.9)(0) = 26+0 = 26. \nonumber
\end{align}
Equity prefers B. Debt then expects $(0.5)(40)+(0.5)(0)=20$, against $40$ under A or C.

\answer{(b)} Strategy B (expected equity payoff \$50 million, against \$35 million for A and
\$26 million for C).

\step{3}{(c) Debt of \$110 million} Same split with $F=110$: equity $\max\{X-110,0\}$ (e.g.\
$\max\{75-110,0\}=0$, $140-110=30$, $300-110=190$), debt $\min\{X,110\}$:
\begin{center}
\begin{tabular}{lccc}
\toprule
Strategy, state (prob.) & Firm & Equity $\max\{X-110,0\}$ & Debt $\min\{X,110\}$\\
\midrule
A (1.0) & 75 & 0 & 75\\
B, high (0.5) & 140 & 30 & 110\\
B, low (0.5) & 0 & 0 & 0\\
C, high (0.1) & 300 & 190 & 110\\
C, low (0.9) & 40 & 0 & 40\\
\bottomrule
\end{tabular}
\end{center}
\begin{align}
  \mathbb{E}[E_A] &= (1.0)(0) = 0, \nonumber\\
  \mathbb{E}[E_B] &= (0.5)(30)+(0.5)(0) = 15+0 = 15, \nonumber\\
  \mathbb{E}[E_C] &= (0.1)(190)+(0.9)(0) = 19+0 = 19. \label{t4e6-EC110}
\end{align}
Equity prefers C. Debt then expects $(0.1)(110)+(0.9)(40)=11+36=47$, against $75$ under A and
$(0.5)(110)+(0.5)(0)=55$ under B.

\answer{(c)} Strategy C (expected equity payoff \$19 million, against \$15 million for B and
\$0 for A).

\step{4}{(d) Expected agency costs} The agency cost is the expected firm payoff of the efficient
strategy minus that of the strategy equity picks. The value-maximizing strategy is A \eqref{t4e6-XA}. With
$F=40$ equity picks B \eqref{t4e6-XB}; with $F=110$ it picks C \eqref{t4e6-XC}:
\begin{align}
  \text{cost}(40) &= \mathbb{E}[X_A]-\mathbb{E}[X_B] = (75)-(70) = 5, \label{t4e6-c40}\\
  \text{cost}(110) &= \mathbb{E}[X_A]-\mathbb{E}[X_C] = (75)-(66) = 9. \label{t4e6-c110}
\end{align}
The expected agency cost is \$5 million with \$40 million of debt and \$9 million with
\$110 million (in payoff-date dollars; divide by $1+r_f$ for today's value).

\answer{(d)} \$5 million with \$40 million of debt; \$9 million with \$110 million of debt (expected
payoff-date values).

\intuition With \$40 million due, A leaves equity a safe \$35 million, but B's 50\% chance of
\$100 million is worth more to equity because the creditors bear the zero-payoff state. With
\$110 million due, A leaves equity nothing and only the long shot C pays off above the face
value often enough. More debt pushes equity towards riskier, lower-value strategies, and the
agency cost grows with leverage. Creditors who anticipate the choice price it in, so the
shareholders bear the cost ex ante.
\pitfalls
\begin{itemize}
  \item Measuring the agency cost as the creditors' loss (e.g.\ $40-20=20$ with $F=40$). That is a
        transfer; the agency cost is the loss in \emph{total} expected payoff, $75-70=5$.
  \item Giving equity $X-F$ in states where $X<F$ (C's low state with $F=110$ gives $-70$, not $0$),
        which makes the risky strategies look worse to equity than they are.
  \item Comparing equity's payoffs across debt levels instead of across strategies at a given debt
        level; each of (b) and (c) is a choice among A, B, C with $F$ fixed.
\end{itemize}
\refbib{Berk Ch.~16, \S~16.6 ``Agency Costs and Information'', ``Equity-Debt Holder Conflicts'' (excessive risk-taking), pp.~483--484. Lecture Topic~4B, slides 2--13.}

% =============================================================================
\exsec{7. Ajax Manufacturing: debt overhang --- Grinblatt \& Titman 16.12--16.13, p.~593}

\asked (16.12) The NPV of the widget project, and whether Ajax can raise the \$75 million when
the existing \$70 million senior bond forbids new debt (so only new equity is possible).
(16.13) Whether Ajax can raise the \$75 million in equity if it uses the riskier, capital-intensive
process instead.

\approach Risk neutrality and $r_f=0$ make every value an expected payoff. New equity is junior to
the bond, so the most new shareholders can be offered is the whole residual claim
$\mathbb{E}[\max\{X-70,0\}]$; the project is fundable iff that is at least \$75 million. This is the
lecture's \emph{debt overhang}: the bond captures part of the project's value. In 16.13 a riskier
payoff shrinks the bond's share (\emph{asset substitution}), which can make the equity issue
feasible.

\notation
\begin{itemize}
  \item $F=70$: existing senior bond obligation, with a covenant forbidding additional debt; $I=75$: investment cost.
  \item States: recession (probability $\pi=0.2$) and favourable (probability $0.8$); $X_s$: project cash flow in state $s$; risk-neutral investors, $r_f=0$, no taxes or bankruptcy costs.
  \item $V$: project value; $E^{\max}$: value of the whole residual claim after the bond, the most new shareholders can be offered; $D$: bond value.
\end{itemize}

\assume The firm has no other assets: without the project, the bond receives nothing. In 16.13
the investment still costs $I=75$ (the statement changes only the cash flows).

\step{1}{(16.12) NPV of the project} We value the project's cash flows and subtract their cost:
\begin{align}
  V &= \pi X_{\text{rec}}+(1-\pi)X_{\text{fav}} \label{t4e7-V}\\
    &= (0.2)(85)+(0.8)(150) \nonumber\\
    &= 17+120 = 137, \nonumber\\
  \mathit{NPV} &= V-I = (137)-(75) = 62. \label{t4e7-NPV}
\end{align}

\step{2}{(16.12) Can it be funded?} We value the largest claim new shareholders can buy and
compare it with the \$75 million they must pay. New debt is forbidden, so the money must come from
equity, which is junior to the \$70 bond. Equity receives $\max\{85-70,0\}=15$ in a recession and
$\max\{150-70,0\}=80$ in the favourable state:
\begin{align}
  E^{\max} &= (0.2)(15)+(0.8)(80) \label{t4e7-E1}\\
           &= 3+64 = 67 < 75 = I. \nonumber
\end{align}
Even 100\% of the equity is worth less than the \$75 million it must raise: \textbf{no}. The
bond, which would be paid in full ($85\ge 70$), would be worth $D=70$ instead of $0$; the
split $70+(67-75)=62$ is the NPV. The creditors capture more than the whole NPV, which is why
the shareholders cannot pay for it.

\answer{(16.12)} $\mathit{NPV}=\$62$ million, but the firm \textbf{cannot} fund the project: all of
the equity is worth \$67 million $<\$75$ million.

\step{3}{(16.13) The capital-intensive variant} We repeat Steps 1--2 with the new cash flows,
$5$ (recession) and $170$ (favourable):
\begin{align}
  V' &= (0.2)(5)+(0.8)(170) \nonumber\\
     &= 1+136 = 137, \nonumber\\
  \mathit{NPV}' &= (137)-(75) = 62, \nonumber
\end{align}
the same NPV but riskier. Equity receives $\max\{5-70,0\}=0$ and $170-70=100$:
\begin{align}
  E'^{\max} &= (0.2)(0)+(0.8)(100) \label{t4e7-E2}\\
            &= 80 \ge 75 = I. \nonumber
\end{align}
\textbf{Yes}: new shareholders pay \$75 million for a fraction $75/80=0.9375$ of the equity and
the old shareholders keep $1-0.9375=6.25\%$, worth $0.0625(80)=80-75=5$. The bond now receives
$\min\{5,70\}=5$ and $\min\{170,70\}=70$, worth
\begin{align}
  D' &= (0.2)(5)+(0.8)(70) \nonumber\\
     &= 1+56 = 57, \nonumber
\end{align}
and $57+80=137=V'$.

\answer{(16.13)} \textbf{Yes}: under the capital-intensive process the equity is worth
\$80 million $\ge\$75$ million, so new shareholders buy $93.75\%$ of it for \$75 million.

\step{4}{Comparison} The two processes create the same value, $62$. The risky one moves
$70-57=13$ from the bondholders to the equity, which is exactly what lifts the equity's
value from $67$ to $80$, above the \$75 million cost. The bondholders prefer the safe process,
but under it nothing is built and they get $0$, so they are better off with the risky one
($57>0$).

\intuition Debt overhang: with a senior bond that cannot be diluted, new equity must pay for
the whole project but receives only the value left after the bond, and here the bond absorbs
too much of it. Raising risk shrinks the bond's share (asset substitution), and that makes the
equity issue possible. Of the two distortions, the risk shift is the lesser here. A cleaner
fix is to renegotiate the bond: under the safe process equity covers \$75 million whenever
the face is cut to at most $137-75=62$, which the bondholders should accept since $62>57$
(the lecture's renegotiation argument).
\pitfalls
\begin{itemize}
  \item Answering ``NPV $>0$, so yes'' in 16.12. Fundability depends on what the junior claim is
        worth to the new money, not on total value.
  \item Valuing equity as $V-F$. That gives $137-70=67$ for \emph{both} processes and the wrong answer
        in 16.13; limited liability makes the recession payoff $\max\{5-70,0\}=0$, not $-65$.
  \item Forgetting that without the project the bond gets nothing here; that is why the bondholders
        prefer the risky process ($57$) to no project ($0$).
\end{itemize}
\refbib{Grinblatt \& Titman Ch.~16, \S~16.2 ``The Debt Overhang Problem'' and ``How the Seniority Structure of Debt Affects the Debt Overhang Problem'', pp.~563--566 (Results~16.3--16.4); ``The Asset Substitution Problem'', pp.~569--572. Lecture Topic~4B, slides 19--31.}

% =============================================================================
\exsec{8. Project riskiness and credit rationing --- Tirole 3.15, p.~150}

\asked Two versions of the fixed-investment project have the same cost, the same expected income
under good behaviour and the same effect of shirking on the success probability, but B succeeds less
often with a larger payoff. Decide which version can be financed for a wider range of the
entrepreneur's cash $A$, i.e.\ is less prone to credit rationing.

\approach In the Holmstr\"om--Tirole fixed-investment model (lecture Topic~4A) a project is
financed iff $A\ge\bar A=I-\mathcal P$, where $\mathcal P=p_H(R-B/\Delta p)$ is the
\emph{pledgeable income}. So the variant with the lower threshold $\bar A$, i.e.\ the higher
pledgeable income, is less rationed. We write $\bar A$ as non-pledgeable rent minus NPV, use equal
NPVs, and compare the rents.

\notation
\begin{itemize}
  \item $j\in\{A,B\}$: project variant; $p_H^j>p_L^j$: success probabilities with good / bad behaviour; $R^j>0$: income on success (0 on failure).
  \item $\Delta p=p_H^A-p_L^A=p_H^B-p_L^B>0$: common; $B>0$: private benefit (common); $I$: common investment; $A$: borrower's cash.
  \item Maintained: $p_H^A R^A=p_H^B R^B$ (equal NPV) and $p_H^A>p_H^B$ (B riskier); (3.1), (3.2) hold for both variants.
  \item $\mathcal P^j=p_H^j(R^j-B/\Delta p)$: pledgeable income; $\bar A^j$: minimum cash for financing variant $j$ (3.3).
\end{itemize}

\step{1}{The minimum stake is the same} We first derive the borrower's minimum reward in each
variant. Variant $j$'s incentive constraint is $p_H^jR_b^j\ge p_L^jR_b^j+B$, i.e.\
$(p_H^j-p_L^j)R_b^j\ge B$, i.e.\ $R_b^j\ge B/\Delta p$: the same for both variants, because $B$ and
$\Delta p$ are common.

\step{2}{Financing thresholds} We compute each variant's threshold from (IR$_l$),
$\mathcal P^j\ge I-A$. From (3.3),
\begin{align}
  \bar A^j &= I-\mathcal P^j \nonumber\\
           &= I-p_H^j\left(R^j-\frac{B}{\Delta p}\right) \nonumber\\
           &= I-p_H^jR^j+p_H^j\frac{B}{\Delta p} \nonumber\\
           &= p_H^j\frac{B}{\Delta p}-(p_H^j R^j-I). \label{t4e8-Abar}
\end{align}
The first term is the non-pledgeable income: the expected rent the borrower must be left.

\step{3}{Compare} Subtracting \eqref{t4e8-Abar} for B from that for A, and substituting
$p_H^B R^B=p_H^A R^A$,
\begin{align}
  \bar A^A-\bar A^B &= \left[p_H^A\frac{B}{\Delta p}-(p_H^A R^A-I)\right]-\left[p_H^B\frac{B}{\Delta p}-((p_H^A R^A)-I)\right] \nonumber\\
                    &= p_H^A\frac{B}{\Delta p}-p_H^AR^A+I-p_H^B\frac{B}{\Delta p}+p_H^AR^A-I \nonumber\\
                    &= (p_H^A-p_H^B)\frac{B}{\Delta p} > 0, \label{t4e8-diff}
\end{align}
positive because $p_H^A>p_H^B$, $B>0$, $\Delta p>0$. The riskier variant B is financed for every
$A\ge\bar A^B$, including $A\in[\bar A^B,\bar A^A)$ where A is not: \textbf{variant B is less prone
to credit rationing}.

\answer{(3.15)} Variant B, the riskier one: $\bar A^B=\bar A^A-(p_H^A-p_H^B)B/\Delta p<\bar A^A$, so
every entrepreneur financed under A is financed under B, and those with
$A\in[\bar A^B,\bar A^A)$ only under B (given that both variants need good behaviour, (3.2)).

\step{4}{Numerical illustration} $I=100$, $\Delta p=0.2$, $B=9$; A: $p_H^A=0.8$, $p_L^A=0.6$,
$R^A=150$; B: $p_H^B=0.6$, $p_L^B=0.4$, $R^B=200$. Both NPVs are $120-100=20$; (3.2) holds
($0.6(150)-100+9=-1<0$ and $0.4(200)-100+9=-11<0$); $B/\Delta p=45$. Then
$\mathcal P^A=0.8(105)=84$, $\bar A^A=16$; $\mathcal P^B=0.6(155)=93$, $\bar A^B=7$; and
$\bar A^A-\bar A^B=(0.2)(45)=9$, as \eqref{t4e8-diff} says.

\step{5}{Beyond the maintained assumption (the book's extension)} If (3.2) fails for the safe
variant, a contract that leaves all income to investors ($R_b=0$, the borrower misbehaves)
may finance it: investors then receive $p_L^AR^A$, and the borrower, who gets the NPV with
shirking, $p_L^AR^A+B-I$, when investors break even, wants it if that NPV is positive. The safe
variant can then be financed and the risky one not when (in order) B cannot be financed with good
behaviour, B cannot be financed with shirking, A can be financed with shirking, and A's NPV with
shirking is positive:
\begin{align}
  I-A &> p_H^B\left(R^B-\frac{B}{\Delta p}\right), \quad I-A>p_L^B R^B, \quad I-A\le p_L^A R^A, \quad I<p_L^A R^A+B. \label{t4e8-ext}
\end{align}
These conditions are compatible. With $A=0$, $I=120$, $\Delta p=0.3$, $B=45$, $p_H^A=0.9$,
$R^A=200$, $p_H^B=0.5$, $R^B=360$ (both $p_H R=180$): $p_H^B(R^B-B/\Delta p)=0.5(210)=105<120$;
$p_L^B R^B=0.2(360)=72<120$; $p_L^A R^A=0.6(200)=120\ge 120$; $p_L^A R^A+B=165>120$. Here the
risky variant is rationed and the safe one is funded. Good behaviour adds more pledgeable
income when the success payoff is large, which favours the risky variant only when good
behaviour is needed for funding.

\intuition The borrower's rent $B/\Delta p$ is paid only on success. A safer project succeeds
more often, so it pays out the rent more often and leaves less expected income to pledge: the
non-pledgeable rent $p_H B/\Delta p$ is higher for the same NPV. Riskier projects with the same mean concentrate
income in success, where it is pledged anyway.
\pitfalls
\begin{itemize}
  \item Importing the Stiglitz--Weiss intuition that riskier borrowers are rationed first. Here
        riskiness lowers the rent $p_HB/\Delta p$ and so \emph{raises} pledgeable income.
  \item Computing the threshold as $I-p_HR$ (all income pledged), which is the same for both variants
        by $p_H^AR^A=p_H^BR^B$ and hides the difference.
  \item Forgetting that the comparison holds only while good behaviour is required for funding: if
        (3.2) fails for one variant, the book's extension (Step 5) can reverse the ranking.
\end{itemize}
\bookans Matches the book's answer (p.~578): with the same minimum reward $B/\Delta p$, the
safe project A is financed for a smaller range of $A$, i.e.\ it is more prone to credit
rationing; the book then shows, with conditions \eqref{t4e8-ext}, that the risky project may be
the more rationed one if high-powered incentives are not necessarily called for. One slip in
the book's prose: it says the third condition means that ``the risky project'' generates
enough pledgeable income when investors get all the cash flow, but that condition,
$I-A\le p_L^A R^A$, is about the safe variant A.
\refbib{Tirole Ch.~3, \S~3.2.1--3.2.2 ``The Fixed-Investment Model'', ``The Lenders' Credit Analysis'', pp.~115--117. Lecture Topic~4A, slides 26--35 (determinants of credit rationing).}

% =============================================================================
\exsec{9. Entrepreneurial risk aversion --- Tirole Exercise 3.2, p.~144}

\asked (i) The entrepreneur is risk neutral above a subsistence level $c_0$ and infinitely averse
to falling below it. Find the smallest cash $\bar A$ with which risk-neutral investors finance the
project, and explain why the answer differs between a project that succeeds for sure under effort
($p_H=1$) and one that can fail even under effort ($p_H<1$). (ii) Redo the analysis for a general
strictly concave utility, once with limited liability (her payments cannot be negative) and once
without.

\approach This is the fixed-investment model of the lecture (Topic~4A) with a risk-averse
borrower. The contract pays her $c^S$ on success and $c^F$ on failure; she chooses it to maximise
expected utility subject to the investors' breakeven constraint (IR$_l$) and her incentive
constraint (IC$_b$), now written in utility. Pledgeable income is total expected income minus the
cheapest expected payment that satisfies (IC$_b$) and the lower bound on $c^F$; a floor on
consumption (subsistence in (i), limited liability in (ii)) raises that payment and so the minimum
equity. In (ii) we write the Lagrangian and derive the first-order conditions, which give the
Holmstr\"om--Shavell insurance-versus-incentives trade-off.

\notation
\begin{itemize}
  \item $A$: entrepreneur's cash ($A\ge c_0$ in (i)); $I>A$: fixed investment; $R>0$: income in
        the case of success, $0$ in the case of failure.
  \item $p_H$, $p_L=p_H-\Delta p$: success probabilities when she works / shirks ($0<p_L<p_H\le 1$);
        $B>0$: private benefit from shirking.
  \item $\mathit{NPV}\equiv p_HR-I>0$; $\bar A_{rn}\equiv I-p_H(R-B/\Delta p)>0$: the
        risk-neutral threshold (3.3), positive by the assumption $I>p_H(R-B/\Delta p)$.
  \item $c^S$, $c^F$: the entrepreneur's consumption in the case of success / failure (her total
        payoff from the contract); $u(\cdot)$: utility of consumption; $c_0$: subsistence level in (i).
  \item $\mathcal{P}$: expected pledgeable income; $\bar A$: minimum cash for financing.
  \item $h(c)\equiv u^{-1}\!\big(u(c)+B/\Delta p\big)$: consumption in success that gives exactly the
        utility premium $B/\Delta p$ over consumption $c$ in failure (part (ii)).
\end{itemize}
\assume The entrepreneur puts all of $A$ into the project and the contract pays her $c^S$ or $c^F$
(keeping cash aside is the same as adding it to both payments). In (ii) the private benefit enters
utility additively, $p\,u(c^S)+(1-p)\,u(c^F)+B$ when she shirks, as in the book's answer.

\step{1}{(i), $p_H<1$: the programme} The entrepreneur offers the contract that maximises her
expected utility subject to the investors' breakeven and her incentive constraint:
\begin{align}
  \max_{c^S,c^F}\; & p_H\,u(c^S)+(1-p_H)\,u(c^F) \label{t4e9-obj}\\
  \text{s.t. }\; & p_H(R-c^S)+(1-p_H)(0-c^F)\ \ge\ I-A, \label{t4e9-irl}\\
  & p_H\,u(c^S)+(1-p_H)\,u(c^F)\ \ge\ p_L\,u(c^S)+(1-p_L)\,u(c^F)+B. \label{t4e9-ic}
\end{align}
Failure has probability $1-p_H>0$ even under work, so $c^F<c_0$ would make \eqref{t4e9-obj} equal to
$-\infty$: the optimum has $c^F\ge c_0$, and then $c^S\ge c_0$ too. On $[c_0,\infty)$ utility is
linear, $u(c)=c$, so \eqref{t4e9-ic} reads
\begin{align}
  p_Hc^S+(1-p_H)c^F &\ge p_Lc^S+(1-p_L)c^F+B \nonumber\\
  (p_H-p_L)c^S-\big[(1-p_L)-(1-p_H)\big]c^F &\ge B \nonumber\\
  (p_H-p_L)c^S-(p_H-p_L)c^F &\ge B \nonumber\\
  (\Delta p)(c^S-c^F) &\ge B. \label{t4e9-iclin}
\end{align}

\step{2}{(i), $p_H<1$: which constraints bind, and $\bar A$} With linear utility the objective
\eqref{t4e9-obj} is the expected payment $p_Hc^S+(1-p_H)c^F$. If \eqref{t4e9-irl} were slack, raising
$c^S$ would raise the objective and relax \eqref{t4e9-iclin}; so \eqref{t4e9-irl} binds and the
entrepreneur gets $A+\mathit{NPV}$. Financing is possible if and only if the cheapest contract that
satisfies \eqref{t4e9-iclin} and $c^F\ge c_0$ leaves the investors enough. Using
$c^S\ge c^F+B/\Delta p$ from \eqref{t4e9-iclin} and then $c^F\ge c_0$, any feasible contract costs
\begin{align}
  p_Hc^S+(1-p_H)c^F &\ge p_H\Big(c^F+\frac{B}{\Delta p}\Big)+(1-p_H)c^F \nonumber\\
                    &= c^F+p_H\frac{B}{\Delta p} \ \ge\ c_0+p_H\frac{B}{\Delta p}, \nonumber
\end{align}
with equality iff both inequalities bind, so the cheapest contract is
\begin{align}
  c^F &= c_0, \label{t4e9-cf}\\
  c^S &= c_0+\frac{B}{\Delta p}. \label{t4e9-cs}
\end{align}
The pledgeable income is total expected income minus this minimum expected payment; substituting
\eqref{t4e9-cf}--\eqref{t4e9-cs},
\begin{align}
  \mathcal{P} &= p_HR-\Big[p_H\Big(c_0+\frac{B}{\Delta p}\Big)+(1-p_H)(c_0)\Big] \nonumber\\
              &= p_H\Big(R-\frac{B}{\Delta p}\Big)-c_0. \label{t4e9-P}
\end{align}
(the middle step: $p_Hc_0+(1-p_H)c_0=c_0$). Financing requires $\mathcal{P}\ge I-A$; using
\eqref{t4e9-P} and moving terms,
\begin{align}
  p_H\Big(R-\frac{B}{\Delta p}\Big)-c_0 &\ge I-A \nonumber\\
  A \ \ge\ \bar A &= I+c_0-p_H\Big(R-\frac{B}{\Delta p}\Big) \label{t4e9-Abar}\\
                  &= \Big[I-p_H\Big(R-\frac{B}{\Delta p}\Big)\Big]+c_0=\bar A_{rn}+c_0. \nonumber
\end{align}
Because $\bar A_{rn}>0$, $\bar A>c_0$: entrepreneurs with $A\in[c_0,\bar A)$ are rationed. The
participation constraint holds, since she gets $A+\mathit{NPV}>A$. Illustration: with
$p_H=0.8$, $p_L=0.4$, $B=1$, $R=10$, $I=6.5$, $c_0=1$, \eqref{t4e9-Abar} gives
$\bar A=(6.5)+(1)-(0.8)(10-1/0.4)=1.5$, against $\bar A_{rn}=0.5$.

\step{3}{(i), $p_H=1$} Now failure never happens on the equilibrium path; it happens only if she
shirks, with probability $1-p_L>0$. Set $c^F<c_0$ (for example $c^F=0$, feasible under limited
liability when $c_0>0$). Shirking then yields $p_L\,u(c^S)+(1-p_L)(-\infty)+B=-\infty$, so
\eqref{t4e9-ic} holds for every $c^S\ge c_0$ and the incentive problem costs nothing. The cheapest
contract is $c^S=c_0$, and
\begin{align}
  \mathcal{P} &= R-c_0, \label{t4e9-P1}\\
  \bar A &= I+c_0-R\ <\ c_0, \label{t4e9-Abar1}
\end{align}
where the inequality uses $R-I=\mathit{NPV}>0$. Every entrepreneur with $A\ge c_0$ is financed: no
rationing (with the numbers above, $\bar A=6.5+1-10=-2.5$). The threshold is discontinuous at
$p_H=1$: by \eqref{t4e9-Abar}, $\bar A\to I+c_0-R+B/\Delta p$ as $p_H\to1$, a jump of $B/\Delta p$.
If $c_0\le0$, limited liability rules out a punishment below $c_0$ and the risk-neutral analysis applies.

\answer{(i)} For $p_H<1$: $\bar A=I+c_0-p_H(R-B/\Delta p)=\bar A_{rn}+c_0$, so entrepreneurs with
$A\in[c_0,\bar A)$ are rationed. For $p_H=1$: failure reveals shirking and can be punished with
$c^F<c_0$ at no cost on the equilibrium path, so $\bar A=I+c_0-R<c_0$ and every entrepreneur with
$A\ge c_0$ is financed (no rationing).

\step{4}{(ii) The general programme} With $u'>0$, $u''<0$, the contract solves
\begin{align}
  \max_{c^S,c^F}\; & p_H\,u(c^S)+(1-p_H)\,u(c^F) \label{t4e9-gobj}\\
  \text{s.t. }\; & p_H(R-c^S)-(1-p_H)c^F\ \ge\ I-A, \tag{IR$_l$}\label{t4e9-girl}\\
  & (\Delta p)\big[u(c^S)-u(c^F)\big]\ \ge\ B, \tag{IC$_b$}\label{t4e9-gic}\\
  & c^F\ge 0,\ c^S\ge0 \quad\text{(only under limited liability)}, \tag{LL}\label{t4e9-ll}\\
  & p_H\,u(c^S)+(1-p_H)\,u(c^F)\ \ge\ u(A). \tag{PC}\label{t4e9-pc}
\end{align}
\eqref{t4e9-gic} is the risk-averse version of (IC$_b$): subtract $p_Lu(c^S)+(1-p_L)u(c^F)+B$
from \eqref{t4e9-gobj}. \eqref{t4e9-pc} says she prefers the project to consuming $A$. In utility
units $v^S=u(c^S)$, $v^F=u(c^F)$ the objective and \eqref{t4e9-gic} are linear and the cost
$p_Hu^{-1}(v^S)+(1-p_H)u^{-1}(v^F)$ is convex, so the problem is convex and the first-order
conditions are sufficient.

\step{5}{(ii) Which constraints bind} We show both (IR$_l$) and (IC$_b$) bind, the second by
deriving the first-order conditions. (a) \eqref{t4e9-girl} binds: otherwise raise $c^S$, which raises
\eqref{t4e9-gobj} and relaxes \eqref{t4e9-gic}. (b) With multipliers
$\lambda\ge0$ on \eqref{t4e9-girl} and $\mu\ge0$ on \eqref{t4e9-gic} (and \eqref{t4e9-ll} slack), the
Lagrangian is
\begin{align}
  \mathcal L &= p_H\,u(c^S)+(1-p_H)\,u(c^F)+\lambda\big[p_H(R-c^S)-(1-p_H)c^F-I+A\big] \nonumber\\
             &\quad+\mu\big[(\Delta p)\big(u(c^S)-u(c^F)\big)-B\big]. \nonumber
\end{align}
Differentiating with respect to $c^S$ and $c^F$ and setting the derivatives to zero,
\begin{align}
  \frac{\partial\mathcal L}{\partial c^S} &= p_Hu'(c^S)-\lambda p_H+\mu(\Delta p)u'(c^S)=0, \nonumber\\
  &\Longrightarrow\quad u'(c^S)\big(p_H+\mu\Delta p\big)=\lambda p_H, \nonumber\\
  \frac{\partial\mathcal L}{\partial c^F} &= (1-p_H)u'(c^F)-\lambda(1-p_H)-\mu(\Delta p)u'(c^F)=0, \nonumber\\
  &\Longrightarrow\quad u'(c^F)\big(1-p_H-\mu\Delta p\big)=\lambda(1-p_H). \nonumber
\end{align}
Here $\lambda>0$, because $\lambda=0$ would make the first condition read $u'(c^S)(p_H+\mu\Delta p)=0$,
impossible with $u'>0$. Dividing the first by $\lambda p_H\,u'(c^S)$ and the second by
$\lambda(1-p_H)\,u'(c^F)$,
\begin{align}
  \frac{1}{u'(c^S)} &= \frac{1}{\lambda}\Big(1+\mu\frac{\Delta p}{p_H}\Big), \label{t4e9-focs}\\
  \frac{1}{u'(c^F)} &= \frac{1}{\lambda}\Big(1-\mu\frac{\Delta p}{1-p_H}\Big). \label{t4e9-focf}
\end{align}
If $\mu=0$, \eqref{t4e9-focs}--\eqref{t4e9-focf} give $u'(c^S)=u'(c^F)$, i.e.\ full insurance
$c^S=c^F$, which violates \eqref{t4e9-gic} because $B>0$. So $\mu>0$ and, by \eqref{t4e9-focs}--\eqref{t4e9-focf},
$c^S>c^F$: the entrepreneur bears risk, but only as much as incentives require
(Holmstr\"om--Shavell inverse-marginal-utility rule).

\step{6}{(ii) The solution} We reduce the two binding constraints to one equation in $c^F$. The
binding \eqref{t4e9-gic} is $u(c^S)=u(c^F)+B/\Delta p$, i.e.\ $c^S=h(c^F)$. The binding
\eqref{t4e9-girl} is $p_HR-p_Hc^S-(1-p_H)c^F=I-A$, i.e.
\begin{align}
  p_Hc^S+(1-p_H)c^F=A+(p_HR-I)=A+\mathit{NPV}. \nonumber
\end{align}
Substituting $c^S=h(c^F)$,
\begin{align}
  p_H\,h(c^F)+(1-p_H)\,c^F = A+\mathit{NPV}. \label{t4e9-sol}
\end{align}
Differentiating the identity $u(h(c))=u(c)+B/\Delta p$ with respect to $c$,
\begin{align}
  u'\big(h(c)\big)\,h'(c)=u'(c) \quad\Longrightarrow\quad h'(c)=\frac{u'(c)}{u'(h(c))}\ \ge\ 1, \nonumber
\end{align}
because $h(c)>c$ and $u'$ is decreasing. Hence the left-hand side of \eqref{t4e9-sol} is increasing
in $c^F$, so \eqref{t4e9-sol} defines a unique $c^F=c(A)$, increasing in $A$. Substituting $u(c^S)=u(c^F)+B/\Delta p$ into
\eqref{t4e9-gobj}, her utility is
\begin{align}
  V(A) &= p_H\Big(u(c(A))+\frac{B}{\Delta p}\Big)+(1-p_H)\,u(c(A)) \nonumber\\
       &= u(c(A))+\frac{p_HB}{\Delta p}. \label{t4e9-V}
\end{align}

\step{7}{(ii) Without limited liability} Only \eqref{t4e9-pc} can fail. Let $\mathit{CE}$ be the certainty
equivalent of the lottery $(c^S,c^F)$, $u(\mathit{CE})=V(A)$, and $\pi\equiv \mathbb{E}[c]-\mathit{CE}$ its
risk premium; by \eqref{t4e9-sol}, $\mathbb{E}[c]=A+\mathit{NPV}$. Then
\begin{align}
  \text{financing}\iff V(A)\ge u(A)\iff \mathit{CE}\ge A\iff \pi(A)\le \mathit{NPV}. \label{t4e9-nll}
\end{align}
A risk-neutral entrepreneur ($\pi=0$) without limited liability is always financed (she can be
punished with $c^F<0$; there is no $\bar A$). A risk-averse one is financed only if the risk premium
of the incentive lottery does not exceed the NPV: risk aversion alone can block positive-NPV projects.
Whether this gives a threshold in $A$ depends on how $\pi$ varies with wealth (with $B$ in utility
units, a richer entrepreneur needs a larger money spread to produce the utility gap $B/\Delta p$).

\step{8}{(ii) With limited liability} \eqref{t4e9-ll} requires $c(A)\ge0$. Because the left-hand side of
\eqref{t4e9-sol} is increasing, $c(A)\ge0$ iff $A+\mathit{NPV}\ge p_Hh(0)$, i.e.
\begin{align}
  A\ \ge\ A_{LL}\equiv I-p_H\big[R-h(0)\big]. \label{t4e9-all}
\end{align}
Below $A_{LL}$ no contract satisfies \eqref{t4e9-girl}, \eqref{t4e9-gic} and \eqref{t4e9-ll}: the cheapest
incentive-compatible contract, $c^F=0$, $c^S=h(0)$, already costs more than $A+\mathit{NPV}$. Above
$A_{LL}$ the solution is that of Step 6 with $c^F=c(A)\ge0$, so the minimum equity is
$\bar A=\min\{A\ge A_{LL}:\ V(A)\ge u(A)\}$. Two remarks. First, for $A>A_{LL}$ the entrepreneur
receives $c^F>0$ in failure: unlike the risk-neutral case, pure debt ($c^F=0$) is no longer optimal,
because she buys partial insurance. Second, risk neutrality gives back the book: with $u(c)=c$,
$h(0)=B/\Delta p$ and \eqref{t4e9-all} is \eqref{t4e9-Abar} with $c_0=0$; part (i) is the case of a floor
$c_0$ instead of $0$. If $u$ is normalised so that $u'(0)=1$, concavity gives $h(0)\ge B/\Delta p$, so
$A_{LL}\ge\bar A_{rn}$: risk aversion raises the equity requirement.

\answer{(ii)} Both (IR$_l$) and (IC$_b$) bind; the contract solves \eqref{t4e9-sol} with
$u(c^S)=u(c^F)+B/\Delta p$, and the first-order conditions \eqref{t4e9-focs}--\eqref{t4e9-focf} give
$c^S>c^F$ (partial insurance). Without limited liability, financing obtains iff the risk premium of
the incentive lottery is at most the NPV, $\pi(A)\le\mathit{NPV}$. With limited liability, also
$A\ge A_{LL}=I-p_H[R-h(0)]$, and $\bar A=\min\{A\ge A_{LL}:V(A)\ge u(A)\}$; above $A_{LL}$ she gets
$c^F>0$, so pure debt is no longer optimal.

\intuition Incentives need a spread between the success and failure payoffs. A risk-neutral
entrepreneur bears that spread for free; a subsistence floor or limited liability bounds the
punishment, so the spread must be paid as a rent in success, and a risk-averse entrepreneur must in
addition be compensated for the risk. Both raise the stake she keeps and lower pledgeable income.
The $p_H=1$ case is the exception: failure then proves shirking, so an unbounded punishment off the
equilibrium path removes the agency cost.

\pitfalls
\begin{itemize}
  \item Subtracting $c_0$ from the cash instead of adding it to the payments: the entrepreneur must be
        paid at least $c_0$ in \emph{both} states, which costs investors $c_0$ in expectation; the two
        bookkeepings agree only if this is done consistently.
  \item Treating $p_H=1$ like $p_H<1$. With $p_H=1$ failure never occurs under effort, so the
        $-\infty$ punishment is free; with $p_H<1$ it would hit her on the equilibrium path.
  \item In (ii), carrying over the risk-neutral result that pure debt ($c^F=0$) is optimal. A
        risk-averse borrower buys partial insurance, $c^F>0$, whenever $A>A_{LL}$.
\end{itemize}

\bookans Matches the book's answer (p.~571): $\bar A=I+c_0-p_H(R-B/\Delta p)$ for $p_H<1$, the
programme of (ii) with binding incentive constraint. For $p_H=1$ the book writes the pledgeable
income as $p_HR$, i.e.\ with $c_0$ kept aside from $A$; this is \eqref{t4e9-P1}--\eqref{t4e9-Abar1}
in other bookkeeping. The book only states the programme for (ii) and cites Holmstr\"om (1979) and
Shavell (1979); Steps 6--8 are this solution's.

\refbib Tirole \S3.2.1--3.2.2, pp.~115--117; answer p.~571.

% =============================================================================
\exsec{10. Maximal incentives principle in the fixed-investment model --- Tirole Exercise 3.18, p.~151}

\asked In the fixed-investment model, failure now yields a positive income $R^F$ (salvage value)
instead of $0$. Show that a debt contract (investors take everything in failure and a fixed amount
in success) is an optimal financial structure, and show that, unlike the variable-investment model
of Section~3.4.3, other contracts may be optimal too.

\approach This is the lecture's ``debt dominates equity'' result (maximal incentives principle,
Topic~4A). Set up the borrower's programme: maximise her expected reward subject to (IC$_b$),
(IR$_l$) and limited liability. With a fixed investment her utility equals the NPV for \emph{every}
contract on which (IR$_l$) binds, so the contract matters only for feasibility, i.e.\ for
\emph{pledgeable income}. Paying her in failure buys no incentive, so the contract that maximises
pledgeable income pays her nothing in failure: debt. Uniqueness fails because, when $A>\bar A$, the
spare pledgeable income can be paid in either state.

\notation
\begin{itemize}
  \item $I$: fixed investment; $A$: entrepreneur's cash; $R^S>R^F>0$: income in success / failure,
        with $R^F<I-A$; $R\equiv R^S-R^F$.
  \item $p_H$, $p_L=p_H-\Delta p$, $B$: as in Section 3.2.
  \item $R_b^S$, $R_b^F\ge0$: entrepreneur's reward in success / failure (limited liability).
  \item $\mathit{NPV}\equiv p_HR^S+(1-p_H)R^F-I=p_HR+R^F-I>0$; $U_b$: entrepreneur's net utility.
  \item $D$: face value of debt; $\bar A$: minimum cash for financing.
\end{itemize}

\step{1}{Programme} The borrower designs the contract $(R_b^S,R_b^F)$; her incentive constraint is
$p_HR_b^S+(1-p_H)R_b^F\ge p_LR_b^S+(1-p_L)R_b^F+B$, which simplifies (collect $R_b^S$ and $R_b^F$
terms, $(1-p_L)-(1-p_H)=\Delta p$) to \eqref{t4e10-ic}:
\begin{align}
  U_b=\max_{R_b^S,R_b^F}\; & p_HR_b^S+(1-p_H)R_b^F-A \label{t4e10-obj}\\
  \text{s.t. }\; & (\Delta p)(R_b^S-R_b^F)\ \ge\ B, \tag{IC$_b$}\label{t4e10-ic}\\
  & p_H(R^S-R_b^S)+(1-p_H)(R^F-R_b^F)\ \ge\ I-A, \tag{IR$_l$}\label{t4e10-irl}\\
  & R_b^S\ge0,\ R_b^F\ge0. \tag{LL}\label{t4e10-ll}
\end{align}

\step{2}{The breakeven constraint binds; $U_b$ is the NPV} If \eqref{t4e10-irl} were slack, raising
$R_b^S$ would raise \eqref{t4e10-obj} and relax \eqref{t4e10-ic}. With \eqref{t4e10-irl} binding,
solve it for her expected reward,
\begin{align}
  p_HR^S-p_HR_b^S+(1-p_H)R^F-(1-p_H)R_b^F &= I-A \nonumber\\
  p_HR_b^S+(1-p_H)R_b^F &= p_HR^S+(1-p_H)R^F-(I-A), \nonumber
\end{align}
and substitute it into \eqref{t4e10-obj}:
\begin{align}
  U_b &= \big[p_HR^S+(1-p_H)R^F-(I-A)\big]-A \nonumber\\
      &= p_HR^S+(1-p_H)R^F-I \nonumber\\
      &= \mathit{NPV}. \label{t4e10-Ub}
\end{align}
Because $I$ is fixed, \eqref{t4e10-Ub} is the same for \emph{every} feasible contract with binding
\eqref{t4e10-irl}. The financial structure matters only for feasibility: it must deliver
\eqref{t4e10-ic} while leaving investors $I-A$.

\step{3}{Maximal pledgeable income and the financing condition} The constraint set is nonempty iff
the largest expected income that can go to investors subject to \eqref{t4e10-ic} and
\eqref{t4e10-ll} covers $I-A$, i.e.\ iff expected income minus the \emph{smallest} expected reward
compatible with \eqref{t4e10-ic} and \eqref{t4e10-ll} is at least $I-A$. Using
$R_b^S\ge R_b^F+B/\Delta p$ from \eqref{t4e10-ic},
\begin{align}
  p_HR_b^S+(1-p_H)R_b^F &\ge p_H\Big(R_b^F+\frac{B}{\Delta p}\Big)+(1-p_H)R_b^F \nonumber\\
                        &= R_b^F+p_H\frac{B}{\Delta p}\ \ge\ p_H\frac{B}{\Delta p}, \nonumber
\end{align}
the last step by $R_b^F\ge0$; both bounds are attained by $R_b^F=0$, $R_b^S=B/\Delta p$. Each unit of
failure reward costs one unit of expected income and buys no incentive. So
\begin{align}
  \mathcal{P} &= p_HR^S+(1-p_H)R^F-p_H\frac{B}{\Delta p} \nonumber\\
              &= p_H(R^S-R^F)+R^F-p_H\frac{B}{\Delta p} \nonumber\\
              &= R^F+p_H\Big(R-\frac{B}{\Delta p}\Big), \label{t4e10-P}
\end{align}
and financing obtains iff $\mathcal{P}\ge I-A$, i.e.
\begin{align}
  A\ \ge\ \bar A\equiv I-R^F-p_H\Big(R-\frac{B}{\Delta p}\Big)=p_H\frac{B}{\Delta p}-\mathit{NPV}, \label{t4e10-Abar}
\end{align}
where the second form uses $\mathit{NPV}=p_HR+R^F-I$.
The safe income $R^F$ works exactly like cash on hand.

\step{4}{A debt contract is optimal} Let investors hold debt of face value $D$ defined by
\begin{align}
  p_HD+(1-p_H)R^F = I-A. \label{t4e10-D}
\end{align}
In failure the firm defaults and investors take $R^F$; in success they get $D$. By $R^F<I-A$,
\eqref{t4e10-D} gives $D>R^F$, so the borrower gets $R_b^F=0$. Substituting \eqref{t4e10-D} into
\eqref{t4e10-Abar} (written as $p_H(R^S-B/\Delta p)+(1-p_H)R^F\ge I-A$),
\begin{align}
  p_H\Big(R^S-\frac{B}{\Delta p}\Big)+(1-p_H)R^F &\ge p_HD+(1-p_H)R^F \nonumber\\
  p_H\Big(R^S-\frac{B}{\Delta p}\Big) &\ge p_HD \nonumber\\
  R^S-\frac{B}{\Delta p} &\ge D \nonumber\\
  R_b^S=R^S-D &\ge \frac{B}{\Delta p}, \label{t4e10-debtic}
\end{align}
(subtracting $(1-p_H)R^F$, then dividing by $p_H>0$, then rearranging)
so \eqref{t4e10-ic} holds, \eqref{t4e10-irl} binds by construction, and by \eqref{t4e10-Ub} debt
attains the maximum. It is optimal whenever financing is feasible.

\step{5}{Debt is not uniquely optimal} The optimal set is every $(R_b^S,R_b^F)\ge0$ with binding
\eqref{t4e10-irl} and \eqref{t4e10-ic}. From the binding \eqref{t4e10-irl},
$R_b^S=[A+\mathit{NPV}-(1-p_H)R_b^F]/p_H$. Substituting this into \eqref{t4e10-ic},
\begin{align}
  (\Delta p)\Big(\frac{A+\mathit{NPV}-(1-p_H)R_b^F}{p_H}-R_b^F\Big) &\ge B \nonumber\\
  (\Delta p)\,\frac{A+\mathit{NPV}-(1-p_H)R_b^F-p_HR_b^F}{p_H} &\ge B \nonumber\\
  \frac{A+\mathit{NPV}-R_b^F}{p_H} &\ge \frac{B}{\Delta p} \nonumber\\
  A+\mathit{NPV}-R_b^F &\ge p_H\frac{B}{\Delta p} \nonumber\\
  R_b^F &\le A+\mathit{NPV}-p_H\frac{B}{\Delta p}=A-\bar A, \label{t4e10-set}
\end{align}
using \eqref{t4e10-Abar} in the last line. So every $R_b^F\in[0,A-\bar A]$ is optimal (add
$R_b^F\le R^F$ if investors' claim must be nonnegative in each state). The contract is unique only
at $A=\bar A$. For $A>\bar A$, a small failure reward (debt plus a little equity, or even all-equity
when its implied $R_b^F$ satisfies \eqref{t4e10-set}) does as well. Illustration: $R^S=12$, $R^F=2$,
$p_H=0.8$, $p_L=0.4$, $B=1.6$, $I=9$, $A=3$. Then $B/\Delta p=1.6/0.4=4$,
$\mathit{NPV}=0.8(12)+0.2(2)-9=9.6+0.4-9=1$, $\bar A=0.8(4)-1=2.2$, and by \eqref{t4e10-D}
$D=(9-3-0.2(2))/0.8=5.6/0.8=7$ (with $2<7\le12-4=8$). The optimal set is
$R_b^F\in[0,\,3-2.2]=[0,0.8]$; at $R_b^F=0.8$, $R_b^S=(3+1-0.2(0.8))/0.8=3.84/0.8=4.8$ and
\eqref{t4e10-ic} binds: $0.4(4.8-0.8)=1.6=B$.

\step{6}{Contrast with variable investment} In Section 3.4.3 the entrepreneur's utility
$(p_HR+R^F-1)I$ rises with $I$, so she expands $I$ until the financing constraint binds, the
analogue of $A=\bar A$. Then the slack in \eqref{t4e10-set} is zero, $R_b^F=0$ is forced and debt is
the unique optimum; any $R_b^F>0$ would waste incentive power and reduce the borrowing capacity (3.16).

\answer{(3.18)} Debt with face value $D$ given by $p_HD+(1-p_H)R^F=I-A$ (so $R^F<D\le R^S-B/\Delta p$)
is optimal whenever the project can be financed, $A\ge\bar A=I-R^F-p_H(R-B/\Delta p)$. It is not
uniquely optimal: every contract with $R_b^F\in[0,A-\bar A]$ and binding (IR$_l$) gives the same
utility $\mathit{NPV}$; the optimum is unique only at $A=\bar A$.

\intuition Paying the entrepreneur in failure costs expected income but weakens incentives, so it is
the worst way to reward her; debt, which gives her nothing in failure, maximises incentives per unit
of pay. With a fixed project, once the project is funded there is no further use for pledgeable
income, so any leftover slack can be paid out in either state.

\pitfalls
\begin{itemize}
  \item Concluding that all-equity is never optimal, as in Section~3.4.3. That needs the investment to
        expand until the financing constraint binds; with a fixed $I$ and $A>\bar A$ it does not bind.
  \item Leaving $R^F$ out of pledgeable income. The salvage value is not subject to moral hazard and
        works exactly like cash on hand: $\bar A$ falls one for one with $R^F$.
  \item Writing (IC$_b$) as $(\Delta p)R_b^S\ge B$, forgetting that a failure reward $R_b^F$ weakens
        incentives: what matters is the spread $R_b^S-R_b^F$.
\end{itemize}

\bookans Matches the book's answer (p.~579), including condition (1) there,
$p_H(R-B/\Delta p)\ge I-(A+R^F)$, the debt implementation with $R^F<D<R^S$ and the
non-uniqueness when (IR$_l$) holds with slack at $R_b^F=0$; \eqref{t4e10-set} makes the size of the
optimal set explicit.

\refbib Tirole \S3.2.2, pp.~116--117; \S3.4.3, pp.~128--129; answer p.~579.

% =============================================================================
\exsec{11. Retained-earnings benefit --- Tirole Exercise 3.11, p.~148}

\asked (i) Find how much one extra unit of retained earnings (the entrepreneur's date-1 profit kept
in the firm) is worth to her at date 2, given that it helps finance a second project whose private
benefit $B^2$ is still random. (ii) Show, by argument or example, that the existence of the second
project can get the first project financed even though no date-2 income can be pledged at date 1.

\approach Two applications of the fixed-investment model (lecture Topic~4A, investment and internal
funds). At date 2, retained earnings play the role of cash $A$ in (IR$_l$): more of them lets the
second project be financed for higher values of $B^2$. Their \emph{shadow value} is the derivative of
the entrepreneur's gross date-2 utility. At date 1, that continuation value is lost when project 1
fails, so it works as an extra reward for effort: the date-1 incentive constraint needs less cash
reward, and date-1 \emph{pledgeable income} rises.

\notation
\begin{itemize}
  \item Project $t\in\{1,2\}$: investment $I^t$, income $R^t$ in success, probabilities $p_H^t$,
        $p_L^t=p_H^t-\Delta p^t$, private benefit $B^t$; $\mathit{NPV}^2\equiv p_H^2R^2-I^2>0$.
  \item $A^1$: entrepreneur's cash at date 1; $R_b^1$: her date-1 reward in success (her retained
        earnings), $0$ in failure.
  \item $B^2\in[\underline B^2,\bar B^2]$: random, c.d.f.\ $F$, density $f$, revealed at the start of date 2.
  \item $\hat B^2(R)$: largest $B^2$ for which project 2 is financed with retained earnings $R$.
  \item $W(R)$: entrepreneur's expected gross date-2 utility with retained earnings $R$; $\mu(R)\equiv W'(R)$:
        shadow value of retained earnings ($v$ in \S3.4.2).
\end{itemize}
\assume No discounting; the entrepreneur's only date-2 wealth is $R_b^1$ (or $0$), and date-2 lenders
cannot be pledged date-1 income.

\step{1}{(i) Date-2 financing given retained earnings $R$} Project 2 is financed iff its pledgeable
income covers the outside funding, (IR$_l$) for project 2:
\begin{align}
  p_H^2\Big(R^2-\frac{B^2}{\Delta p^2}\Big) &\ge I^2-R. \label{t4e11-f2}
\end{align}
Solving \eqref{t4e11-f2} for $B^2$ (expand, move $p_H^2B^2/\Delta p^2$ to the right and $I^2-R$ to
the left, then multiply by $\Delta p^2/p_H^2>0$),
\begin{align}
  p_H^2R^2-\frac{p_H^2B^2}{\Delta p^2} &\ge I^2-R \nonumber\\
  \frac{p_H^2B^2}{\Delta p^2} &\le p_H^2R^2-I^2+R=\mathit{NPV}^2+R \nonumber\\
  B^2\ \le\ \hat B^2(R) &\equiv \frac{\Delta p^2}{p_H^2}\big(\mathit{NPV}^2+R\big). \label{t4e11-Bhat}
\end{align}
By the assumption $\underline B^2>\Delta p^2\,\mathit{NPV}^2/p_H^2$, \eqref{t4e11-Bhat} gives
$\hat B^2(0)<\underline B^2$: with no retained earnings, project 2 is never financed.

\step{2}{(i) Gross date-2 utility} We follow the hint and compute her expected gross utility. If
financed, project 2's lenders break even and the entrepreneur gets her cash plus $\mathit{NPV}^2$
(equation (3.6)); otherwise she keeps $R$. With probability $F(\hat B^2(R))$ it is financed, so
\begin{align}
  W(R) &= F\big(\hat B^2(R)\big)\big(R+\mathit{NPV}^2\big)+\Big[1-F\big(\hat B^2(R)\big)\Big]R \nonumber\\
       &= R+F\big(\hat B^2(R)\big)\,\mathit{NPV}^2. \label{t4e11-W}
\end{align}

\step{3}{(i) Shadow value} Differentiating \eqref{t4e11-W} and using
$d\hat B^2/dR=\Delta p^2/p_H^2$ from \eqref{t4e11-Bhat},
\begin{align}
  \mu(R)=W'(R) &= 1+f\big(\hat B^2(R)\big)\,\frac{d\hat B^2}{dR}\,\mathit{NPV}^2 \nonumber\\
               &= 1+f\big(\hat B^2(R)\big)\,\frac{\Delta p^2}{p_H^2}\,\mathit{NPV}^2\ \ge 1. \label{t4e11-mu}
\end{align}
One unit of retained earnings is worth one unit of consumption plus the expected NPV it unlocks: it
raises the probability of financing at the margin by $f\,\Delta p^2/p_H^2$.

\answer{(i)} $\mu(R)=1+f\big(\hat B^2(R)\big)\,(\Delta p^2/p_H^2)\,(p_H^2R^2-I^2)\ \ge1$, with
$\hat B^2(R)=(\Delta p^2/p_H^2)\big(p_H^2R^2-I^2+R\big)$ and gross utility
$W(R)=R+F(\hat B^2(R))(p_H^2R^2-I^2)$.

\step{4}{(ii) Date-1 incentive constraint} By \eqref{t4e11-W} and Step 1, $W(0)=0+F(\hat B^2(0))\,\mathit{NPV}^2=0$, so
the entrepreneur works at date 1 iff
\begin{align}
  p_H^1W(R_b^1)+(1-p_H^1)W(0) &\ge p_L^1W(R_b^1)+(1-p_L^1)W(0)+B^1 \nonumber\\
  (p_H^1-p_L^1)W(R_b^1)-(p_H^1-p_L^1)W(0) &\ge B^1 \nonumber\\
  (\Delta p^1)\,W(R_b^1) &\ge B^1 \qquad (W(0)=0) \nonumber\\
  (\Delta p^1)\big[R_b^1+F\big(\hat B^2(R_b^1)\big)\mathit{NPV}^2\big] &\ge B^1, \tag{IC$^1$}\label{t4e11-ic}
\end{align}
the last line substituting \eqref{t4e11-W}.
Let $\tilde R_b^1$ be the smallest $R_b^1$ satisfying \eqref{t4e11-ic}. Whenever
$F(\hat B^2(\tilde R_b^1))>0$,
\begin{align}
  \tilde R_b^1 = \frac{B^1}{\Delta p^1}-F\big(\hat B^2(\tilde R_b^1)\big)\mathit{NPV}^2 < \frac{B^1}{\Delta p^1}. \label{t4e11-Rt}
\end{align}

\step{5}{(ii) Date-1 financing} The entrepreneur maximises $p_H^1W(R_b^1)-A^1$ subject to
\eqref{t4e11-ic} and $p_H^1(R^1-R_b^1)\ge I^1-A^1$. Since $W$ is increasing \eqref{t4e11-mu}, the
breakeven constraint binds, and financing is possible iff
\begin{align}
  p_H^1\big(R^1-\tilde R_b^1\big)\ \ge\ I^1-A^1. \label{t4e11-fin}
\end{align}
By \eqref{t4e11-Rt} the left-hand side exceeds the stand-alone pledgeable income
$p_H^1(R^1-B^1/\Delta p^1)$. So for $A^1$ with
$p_H^1(R^1-B^1/\Delta p^1)<I^1-A^1\le p_H^1(R^1-\tilde R_b^1)$, project 1 is funded only
because project 2 exists, although none of project 2's income is pledged at date 1.

\step{6}{(ii) An example} Project 2: $I^2=20$, $R^2=40$, $p_H^2=0.75$, $p_L^2=0.25$, $B^2\sim U[7,9]$,
so $\mathit{NPV}^2=10$ and $\Delta p^2\mathit{NPV}^2/p_H^2=20/3<7$ (assumption holds; also $p_L^2R^2+B^2<I^2$).
By \eqref{t4e11-Bhat}, $\hat B^2(R)=\tfrac23(10+R)$, which lies in $[7,9]$ for $R\in[0.5,3.5]$; there
$F=(\hat B^2-7)/2$ and \eqref{t4e11-W}--\eqref{t4e11-mu} give
\begin{align}
  W(R) &= R+\frac{\tfrac23(10+R)-7}{2}\,(10) \nonumber\\
       &= R+\frac{(20+2R-21)/3}{2}\,(10) = R+\frac{5(2R-1)}{3} \nonumber\\
       &= \frac{3R+10R-5}{3}=\frac{13R-5}{3}, \nonumber\\
  \mu &= 1+\Big(\frac12\Big)\Big(\frac{0.5}{0.75}\Big)(10)=1+\frac{10}{3}=\frac{13}{3}. \nonumber
\end{align}
Project 1: $I^1=4.5$, $R^1=10$, $p_H^1=0.5$, $p_L^1=0.25$, $B^1=0.5$, so $B^1/\Delta p^1=2$. Alone, the
pledgeable income is $(0.5)(10-2)=4$ and $\bar A^1=4.5-4=0.5$. With project 2, \eqref{t4e11-ic} at
equality reads $W(\tilde R_b^1)=B^1/\Delta p^1$, i.e.\ $(13\tilde R_b^1-5)/3=2$, so
$13\tilde R_b^1=6+5=11$, $\tilde R_b^1=11/13\approx0.846154$ (inside $[0.5,3.5]$, where the formula
for $W$ holds), and
\begin{align}
  p_H^1(R^1-\tilde R_b^1) = (0.5)\Big(10-\frac{11}{13}\Big)=\frac{119}{26}\approx 4.576923 > 4.5. \nonumber
\end{align}
The threshold falls from $0.5$ to $4.5-119/26=-1/13\approx-0.076923$: any $A^1\in[0,0.5)$ is now funded
(at $A^1=0$, $R_b^1=10-4.5/0.5=1\ge\tilde R_b^1$).

\answer{(ii)} Yes. Because project 2's rent is lost on failure, the date-1 reward needed for effort
falls to $\tilde R_b^1<B^1/\Delta p^1$ and date-1 pledgeable income rises to
$p_H^1(R^1-\tilde R_b^1)>p_H^1(R^1-B^1/\Delta p^1)$. Any $A^1$ with
$p_H^1(R^1-B^1/\Delta p^1)<I^1-A^1\le p_H^1(R^1-\tilde R_b^1)$ is financed only thanks to project 2
(in the example, $A^1\in[0,0.5)$).

\intuition Retained earnings are worth more than their face value because they are the net worth that
gets the next project financed. That future rent is lost if project 1 fails, so it acts as a reward for
date-1 effort: part of the incentive is paid with continuation value instead of cash, and the cash
saved becomes pledgeable to date-1 lenders.

\pitfalls
\begin{itemize}
  \item Valuing retained earnings at face ($\mu=1$). Each unit also raises the chance that project 2 is
        financed, worth $f\,(\Delta p^2/p_H^2)\,\mathit{NPV}^2$ more.
  \item Writing the date-1 incentive constraint with $R_b^1$ alone, $(\Delta p^1)R_b^1\ge B^1$: the
        reward for success is $W(R_b^1)$, which includes the continuation rent.
  \item Treating project 2's income as pledgeable at date 1. The effect works only through the
        entrepreneur's incentives; date-1 lenders are repaid out of $R^1$.
\end{itemize}

\bookans Matches the book's answer (p.~575): $\mu=1+f(\hat B^2)(d\hat B^2/dR_b^1)\mathit{NPV}^2$ and the
date-1 constraint $(\Delta p^1)[R_b^1+F(\hat B^2(R_b^1))\mathit{NPV}^2]\ge B^1$. The book writes the
constraint without $W(0)$; this is right because the assumption on $\underline B^2$ makes $W(0)=0$
(Step 1). The numerical example is this solution's.

\refbib Tirole \S3.2.2, pp.~116--117; \S3.4.2, p.~128 (shadow value of equity); answer p.~575.

% =============================================================================
\exsec{12. Renegotiation and debt forgiveness --- Tirole Exercise 3.6, p.~146}

\asked In the variable-investment model, suppose the initial loan leaves the borrower too small a
stake to work, $R_b<BI/\Delta p$, but can be renegotiated after the investment and before effort.
(i) Show that a mutually beneficial renegotiation exists exactly when the stated inequality holds.
(ii) Explain that inequality, and derive it as the financing condition of a fictitious
fixed-investment project. (iii) If the borrower makes a take-it-or-leave-it renegotiation offer,
find how much she can invest (her borrowing capacity) with such an initial contract. (iv) Argue
directly that writing $R_b\ge BI/\Delta p$ from the start loses nothing.

\approach Lecture Topic~4B, renegotiation. Without renegotiation the borrower shirks; renegotiation
can only help if it restores (IC$_b$), and it happens iff both sides gain relative to the shirking
status quo. The lender gives up part of his claim (\emph{debt forgiveness}) and is paid back out of
the extra \emph{pledgeable income} that effort creates, so the condition has the form ``pledgeable
income $\ge$ investment $-$ cash'' of the fixed-investment model. For (iii) the lender's ex ante
breakeven uses his renegotiated payoff; for (iv) rational expectations make the renegotiated contract
one that could have been signed directly.

\notation
\begin{itemize}
  \item $I$: (sunk) investment; income $RI$ in success, $0$ in failure; $A$: borrower's cash.
  \item $p_H$, $p_L=p_H-\Delta p$; private benefit $BI$ from shirking; (3.7) $p_HR>1$, (3.8)
        $1>p_LR+B$, (3.9) $p_HR<1+p_HB/\Delta p$; these imply $(\Delta p)R>B$.
  \item $R_b<BI/\Delta p$: borrower's stake in the initial contract; $\hat R_b$: renegotiated stake.
  \item $k=1/[1-p_H(R-B/\Delta p)]$: equity multiplier (3.12).
\end{itemize}

\step{1}{(i) Status quo} We compute each party's payoff if nothing is renegotiated; these are the
outside options in the renegotiation. With $R_b<BI/\Delta p$, $(\Delta p)R_b<BI$: (IC$_b$) fails and
the borrower shirks (success probability $p_L$, private benefit $BI$). Payoffs if the contract is not
renegotiated:
\begin{align}
  U_b^0 &= p_LR_b+BI, \label{t4e12-ub0}\\
  U_l^0 &= p_L(RI-R_b). \label{t4e12-ul0}
\end{align}

\step{2}{(i) Renegotiation must restore incentives} A new stake $\hat R_b<BI/\Delta p$ still induces
shirking, so total surplus $p_LRI+BI$ is unchanged and the change is a pure transfer: one side loses.
A mutually advantageous renegotiation therefore needs $\hat R_b\ge BI/\Delta p$, and both parties must
gain relative to \eqref{t4e12-ub0}--\eqref{t4e12-ul0}:
\begin{align}
  p_H\hat R_b &\ge p_LR_b+BI, \label{t4e12-pb}\\
  p_H(RI-\hat R_b) &\ge p_L(RI-R_b). \label{t4e12-pl}
\end{align}

\step{3}{(i) The condition} For any $\hat R_b\ge BI/\Delta p$, \eqref{t4e12-pb} holds automatically, since
$R_b<BI/\Delta p$:
\begin{align}
  p_H\hat R_b \ge p_H\frac{BI}{\Delta p}=BI+p_L\frac{BI}{\Delta p} > BI+p_LR_b. \nonumber
\end{align}
The lender prefers the smallest admissible stake, $\hat R_b=BI/\Delta p$. So a mutually advantageous
$\hat R_b$ exists iff \eqref{t4e12-pl} holds at $\hat R_b=BI/\Delta p$:
\begin{align}
  p_H\Big(RI-\frac{BI}{\Delta p}\Big) &\ge p_L(RI-R_b) \nonumber\\
  p_HRI-\frac{p_HBI}{\Delta p} &\ge p_LRI-p_LR_b \nonumber\\
  (p_H-p_L)RI-\frac{p_HBI}{\Delta p}+p_LR_b &\ge 0 \nonumber\\
  (\Delta p)RI-\frac{p_HBI}{\Delta p}+p_LR_b &\ge 0. \label{t4e12-cond}
\end{align}
Necessity: any renegotiated contract has $\hat R_b\ge BI/\Delta p$, and \eqref{t4e12-pl} is hardest to
meet at larger $\hat R_b$. Sufficiency: if \eqref{t4e12-cond} holds, $\hat R_b=BI/\Delta p$ makes the
lender weakly and the borrower strictly better off. The renegotiation is \emph{debt forgiveness}: the
lender's claim in success falls from $RI-R_b$ to $RI-\hat R_b$.

\answer{(i)} Renegotiation (to $\hat R_b\ge BI/\Delta p$) takes place iff
$(\Delta p)RI-p_HBI/\Delta p+p_LR_b\ge0$.

\step{4}{(ii) Interpretation via the general theory} Treat ``inducing effort'' as a fictitious
fixed-investment project. Its income is the increase in expected output, $(\Delta p)RI$; it needs no
new investment ($I_f=0$); the borrower's minimum stake in it is $p_H\,BI/\Delta p$ (her expected rent
under (IC$_b$)), and she brings as ``cash'' the expected income she gives up, $A_f\equiv p_LR_b$. Its
pledgeable income is $(\Delta p)RI-p_HBI/\Delta p$, and the financing condition (IR$_l$),
``pledgeable income $\ge$ investment $-$ cash'', reads
\begin{align}
  (\Delta p)RI-\frac{p_HBI}{\Delta p} \ \ge\ 0-p_LR_b, \label{t4e12-fict}
\end{align}
which is \eqref{t4e12-cond}. The total gain from restoring effort, $(\Delta p)RI-BI$, is always positive,
yet renegotiation may fail: the lender can only be paid out of pledgeable income, and the borrower's
initial stake $p_LR_b$ plays the role of net worth. A larger $R_b$ makes forgiveness easier.

\answer{(ii)} The condition says the pledgeable part of the effort gain, $(\Delta p)RI-p_HBI/\Delta p$,
must cover the lender's ``outlay'' $0-p_LR_b$: it is (IR$_l$) for a fictitious project with income
$(\Delta p)RI$, investment $0$ and cash $p_LR_b$.

\step{5}{(iii) Borrowing capacity under renegotiation, borrower has all bargaining power} When
\eqref{t4e12-cond} holds, the borrower offers the $\hat R_b$ that leaves the lender indifferent:
\begin{align}
  p_H(RI-\hat R_b) = p_L(RI-R_b), \quad\text{i.e.}\quad \hat R_b=RI-\frac{p_L}{p_H}(RI-R_b). \label{t4e12-tioli}
\end{align}
The lender therefore anticipates exactly \eqref{t4e12-ul0}, and his ex ante breakeven constraint is
\begin{align}
  p_L(RI-R_b) \ \ge\ I-A. \label{t4e12-irl}
\end{align}
The borrower's utility is total surplus with effort minus the lender's payoff:
$p_HRI-p_L(RI-R_b)-A$, which by the binding \eqref{t4e12-irl} equals $p_HRI-(I-A)-A=(p_HR-1)I$, the
NPV with effort. Since $p_HR>1$ (3.7), she maximises $I$ subject to
\eqref{t4e12-cond}, \eqref{t4e12-irl} and $0\le R_b<BI/\Delta p$. \eqref{t4e12-cond} bounds $p_LR_b$
from below and \eqref{t4e12-irl} (expanded: $p_LRI-p_LR_b\ge I-A$) bounds it from above:
\begin{align}
  \frac{p_HBI}{\Delta p}-(\Delta p)RI\ \le\ p_LR_b\ \le\ p_LRI-(I-A). \label{t4e12-band}
\end{align}
An $R_b$ exists iff the lower bound is below the upper bound:
\begin{align}
  \frac{p_HBI}{\Delta p}-(\Delta p)RI &\le p_LRI-I+A \nonumber\\
  I-(\Delta p)RI-p_LRI+\frac{p_HBI}{\Delta p} &\le A \nonumber\\
  I-p_HRI+\frac{p_HBI}{\Delta p} &\le A \qquad (\Delta p+p_L=p_H) \nonumber\\
  I\Big[1-p_H\Big(R-\frac{B}{\Delta p}\Big)\Big] &\le A \nonumber\\
  I &\le kA. \label{t4e12-kA}
\end{align}
If $p_HB/\Delta p\ge(\Delta p)R$, the lower bound in \eqref{t4e12-band} is nonnegative and the capacity
is exactly $kA$: at $I=kA$ both \eqref{t4e12-cond} and \eqref{t4e12-irl} bind, $R_b<BI/\Delta p$ (by
$(\Delta p)R>B$) and \eqref{t4e12-tioli} gives $\hat R_b=BI/\Delta p$. If $p_HB/\Delta p<(\Delta p)R$, the
constraint $R_b\ge0$ binds first, and \eqref{t4e12-irl} with $R_b=0$ gives $I\le A/(1-p_LR)<kA$.
Either way the capacity is at most $kA$, the capacity without renegotiation.

\answer{(iii)} The borrowing capacity is $I=kA$, $k=1/[1-p_H(R-B/\Delta p)]$, when
$p_HB/\Delta p\ge(\Delta p)R$ (then $p_LR_b=kA[p_HB/\Delta p-(\Delta p)R]$ and $\hat R_b=BI/\Delta p$),
and $I=A/(1-p_LR)<kA$ otherwise; it never exceeds the capacity $kA$ of an incentive-compatible
initial contract.

\step{6}{(iv) Rational expectations} Take any initial contract $(I,R_b)$ with $R_b<BI/\Delta p$. If it is
not renegotiated, the borrower shirks; the NPV is $(p_LR+B-1)I<0$ by (3.8), and since the lender
breaks even the borrower's utility is that negative NPV: worse than not investing. If it is
renegotiated to $\hat R_b\ge BI/\Delta p$, the final contract $(I,\hat R_b)$ satisfies (IC$_b$), and by
rational expectations the lender's ex ante breakeven uses his final payoff,
$p_H(RI-\hat R_b)\ge I-A$, which is (IR$_l$). So $(I,\hat R_b)$ could have been signed directly, with the
same effort and the same payoffs. Nothing is lost by restricting attention to $R_b\ge BI/\Delta p$,
which is renegotiation-proof.

\answer{(iv)} Any initial contract with $R_b<BI/\Delta p$ either is not renegotiated (negative NPV,
worse than not investing) or ends as an incentive-compatible contract that satisfies (IR$_l$) and
could have been signed at the outset; so restricting to $R_b\ge BI/\Delta p$ loses nothing.

\intuition A lender facing a borrower who is about to shirk is willing to forgive debt, because the
higher stake raises the success probability and the lender keeps part of that gain. He forgives only
if his remaining claim after forgiveness, the pledgeable part of the effort gain, beats what he gets
from a shirking borrower. Anticipating this, an initial contract that relies on forgiveness can at
best reproduce the incentive-compatible contract: the multiplier $k$ is the right one.

\pitfalls
\begin{itemize}
  \item Arguing that renegotiation always happens because restoring effort raises total surplus
        ($(\Delta p)RI>BI$). The lender gains only the \emph{pledgeable} part of that surplus.
  \item Writing the lender's ex ante breakeven in (iii) as $p_H(RI-R_b)\ge I-A$, i.e.\ his initial
        claim with effort. He anticipates the renegotiation and values his final payoff, which with a
        take-it-or-leave-it offer by the borrower is $p_L(RI-R_b)$.
  \item Mixing the fixed- and variable-investment notation: here the private benefit is $BI$ and the
        minimum stake is $BI/\Delta p$, not $B/\Delta p$.
\end{itemize}

\bookans Matches the book's answer to (i)--(ii) (pp.~573--574), including the fictitious-project reading
with cash $\hat A=p_LR_b$. The book gives no answer to (iii)--(iv); Steps 5--6 are this solution's.

\refbib Tirole \S3.4.1--3.4.2, pp.~127--128; \S3.2.2, pp.~116--117; answer pp.~573--574.

% =============================================================================
\exsec{13. Loan size and collateral requirements --- Tirole Exercise 4.17, p.~195}

\asked The project has decreasing returns, $R(I)$ concave, and the entrepreneur can pledge costly
collateral $C$ that investors seize on failure. Assuming the first best cannot be financed, show
that a poorer entrepreneur (lower $A$) or a worse agency problem (higher $B$, or higher $p_H/\Delta p$
with $p_H$ fixed) leads to a smaller optimal investment and to more collateral.

\approach This combines the variable-investment model (\S3.4) with collateral pledging (\S4.3.5);
lecture Topic~4A, collateral reading. Collateral raises \emph{pledgeable income} twice: investors get
$\phi(C)$ on failure, and the borrower's stake grows, so her cash reward can shrink by $C$ ($p_HC$).
It costs a deadweight loss $C-\phi(C)$ on failure. We eliminate $R_b$ with the binding (IC$_b$),
write the Lagrangian with multiplier $\mu$ on (IR$_l$), derive the first-order conditions, and do
comparative statics on the resulting system with the implicit-function theorem.

\notation
\begin{itemize}
  \item $I\ge0$: investment; $R(I)$: income in success ($R(0)=0$, $R'>0$, $R''<0$, Inada).
  \item $p_H$, $p_L=p_H-\Delta p$; private benefit $BI$; $b\equiv p_HB/\Delta p$: per-unit agency cost
        (rises with $B$ and, with $p_H$ fixed, with $p_H/\Delta p$).
  \item $C\ge0$: collateral pledged in failure, cost $C$ to the entrepreneur, value $\phi(C)$ to
        investors ($\phi(0)=0$, $\phi'>0$, $\phi''<0$, $\phi'(0)=1$, $\phi'(\infty)=0$).
  \item $R_b\ge0$: reward in success; $A$: cash; $I^*$: first-best size, $p_HR'(I^*)=1$.
  \item $\mu\ge0$: multiplier of the breakeven constraint; $\theta\equiv\mu/(1+\mu)\in[0,1)$;
        $\varepsilon(I)\equiv -IR''(I)/R'(I)>0$: curvature of $R$.
\end{itemize}
\assume The limited-liability constraint $R_b\ge0$, i.e.\ $C\le BI/\Delta p$, does not bind (as in
footnote 24 of \S4.3.5).

\step{1}{Programme} We set up the problem and eliminate $R_b$. The entrepreneur gets $R_b$ on
success and $-C$ on failure, so (IC$_b$) is
$p_HR_b-(1-p_H)C\ge p_LR_b-(1-p_L)C+BI$, i.e.\ $(\Delta p)R_b+(\Delta p)C\ge BI$, i.e.\
$(\Delta p)(R_b+C)\ge BI$. Investors get $R(I)-R_b$ in success and $\phi(C)$ in failure. Collateral
destroys $C-\phi(C)\ge0$ in failure, so with binding breakeven her utility is the NPV:
\begin{align}
  \max_{I,C,R_b}\; & U_b=p_HR(I)-I-(1-p_H)\big[C-\phi(C)\big] \label{t4e13-obj}\\
  \text{s.t. }\; & (\Delta p)(R_b+C)\ \ge\ BI, \tag{IC$_b$}\label{t4e13-ic}\\
  & p_H\big[R(I)-R_b\big]+(1-p_H)\phi(C)\ \ge\ I-A. \tag{IR$_l$}\label{t4e13-irl}
\end{align}
$R_b$ does not enter \eqref{t4e13-obj} and lowers \eqref{t4e13-irl}, so \eqref{t4e13-ic} binds:
$R_b=BI/\Delta p-C$. Substituting into \eqref{t4e13-irl},
\begin{align}
  p_H\Big[R(I)-\frac{BI}{\Delta p}+C\Big]+(1-p_H)\phi(C) &\ge I-A \nonumber\\
  p_HR(I)-\frac{p_HB}{\Delta p}\,I+p_HC+(1-p_H)\phi(C)-I+A &\ge 0 \nonumber\\
  G(I,C)\equiv p_HR(I)-bI+p_HC+(1-p_H)\phi(C)-I+A\ &\ge\ 0. \label{t4e13-G}
\end{align}
Collateral raises pledgeable income directly ($\phi$) and by replacing reward ($p_HC$).
The objective is concave and $G$ is concave in $(I,C)$, so the Kuhn--Tucker conditions are sufficient.

\step{2}{First-order conditions} We derive the conditions that characterise the optimum. With
multiplier $\mu\ge0$ on \eqref{t4e13-G}, the Lagrangian is
\begin{align}
  \mathcal L(I,C,\mu) &= p_HR(I)-I-(1-p_H)\big[C-\phi(C)\big] \nonumber\\
     &\quad+\mu\big[p_HR(I)-bI+p_HC+(1-p_H)\phi(C)-I+A\big]. \nonumber
\end{align}
Differentiating with respect to $I$ and $C$ and setting the derivatives to zero,
\begin{align}
  \frac{\partial\mathcal L}{\partial I}&=p_HR'(I)-1+\mu\big[p_HR'(I)-b-1\big]=0, \nonumber\\
  \frac{\partial\mathcal L}{\partial C}&=-(1-p_H)\big[1-\phi'(C)\big]+\mu\big[p_H+(1-p_H)\phi'(C)\big]=0. \nonumber
\end{align}
\emph{The $I$-condition.} Collect the terms in $p_HR'(I)-1$:
\begin{align}
  (1+\mu)\big[p_HR'(I)-1\big]-\mu b &= 0 \nonumber\\
  p_HR'(I)-1 &= \frac{\mu}{1+\mu}\,b. \nonumber
\end{align}
\emph{The $C$-condition.} Expand and collect the terms in $\phi'(C)$:
\begin{align}
  -(1-p_H)+(1-p_H)\phi'(C)+\mu p_H+\mu(1-p_H)\phi'(C) &= 0 \nonumber\\
  (1+\mu)(1-p_H)\phi'(C) &= (1-p_H)-\mu p_H \nonumber\\
  \phi'(C) &= \frac{1}{1+\mu}-\frac{\mu}{1+\mu}\,\frac{p_H}{1-p_H}. \nonumber
\end{align}
Writing $\theta=\mu/(1+\mu)$, so that $1/(1+\mu)=1-\theta$,
\begin{align}
  p_HR'(I)-1 &= \theta\,b, \label{t4e13-focI}\\
  \phi'(C) &= 1-\theta-\theta\frac{p_H}{1-p_H}=1-\theta\,\frac{(1-p_H)+p_H}{1-p_H}=1-\frac{\theta}{1-p_H}. \label{t4e13-focC}
\end{align}
If $\mu=0$, \eqref{t4e13-focI}--\eqref{t4e13-focC} give $I=I^*$, $C=0$, which violates \eqref{t4e13-G}
by the assumption $p_H[R(I^*)-BI^*/\Delta p]<I^*-A$. So $\theta>0$: \eqref{t4e13-focI} gives $I<I^*$ and
\eqref{t4e13-focC} gives $\phi'(C)<1=\phi'(0)$, i.e.\ $C>0$; finiteness of $C$ needs
$\theta<1-p_H$. These are the book's conditions, since $\theta=\mu/(1+\mu)$.

\step{3}{Monotonicity in $\theta$} The FOCs define $I$ and $C$ as functions of $\theta$ (and $b$);
we find how they move and then pin $\theta$ down with the binding breakeven. Totally differentiating
\eqref{t4e13-focI}, $p_HR''(I)\,dI=b\,d\theta+\theta\,db$; and \eqref{t4e13-focC},
$\phi''(C)\,dC=-d\theta/(1-p_H)$. With $R''<0$, $\phi''<0$ and $Q\equiv1/(p_H|R''(I)|)>0$, dividing by
$p_HR''=-1/Q$ and by $\phi''$:
\begin{align}
  \frac{\partial I}{\partial\theta} &= -bQ<0,\qquad \frac{\partial I}{\partial b}=-\theta Q<0, \label{t4e13-Ith}\\
  C'(\theta) &= -\frac{1}{(1-p_H)\phi''(C)}>0. \label{t4e13-Cth}
\end{align}
Let $g(\theta;A,b)\equiv G(I(\theta,b),C(\theta))$, the slack of \eqref{t4e13-G} along the FOCs; the
optimum is the root $g=0$. Using \eqref{t4e13-focI}, $p_HR'=1+\theta b$, so
$G_I=p_HR'-b-1=(1+\theta b)-b-1=-(1-\theta)b<0$; also $G_C=p_H+(1-p_H)\phi'>0$. Then, by the chain
rule and \eqref{t4e13-Ith}--\eqref{t4e13-Cth},
\begin{align}
  g_\theta &= G_I\frac{\partial I}{\partial\theta}+G_CC'(\theta) \nonumber\\
           &= \big(-(1-\theta)b\big)\big(-bQ\big)+G_CC'(\theta) \nonumber\\
           &= (1-\theta)b^2Q+G_CC'(\theta)\ >0. \label{t4e13-gth}
\end{align}
$g$ is $<0$ at $\theta=0$ (first best not pledgeable) and $\to+\infty$ as $\theta\to1-p_H$ ($C\to\infty$),
so the root $\theta^*\in(0,1-p_H)$ is unique.

\step{4}{Comparative statics in $A$} $g_A=1$, so by the implicit-function theorem and \eqref{t4e13-gth}
\begin{align}
  \frac{d\theta^*}{dA} = -\frac{g_A}{g_\theta} = -\frac{1}{g_\theta}<0. \label{t4e13-thA}
\end{align}
Combining with \eqref{t4e13-Ith}--\eqref{t4e13-Cth} by the chain rule:
\begin{align}
  \frac{dI}{dA} &= \frac{\partial I}{\partial\theta}\,\frac{d\theta^*}{dA}=(-bQ)\Big(-\frac{1}{g_\theta}\Big)=\frac{bQ}{g_\theta}>0, \nonumber\\
  \frac{dC}{dA} &= C'(\theta)\,\frac{d\theta^*}{dA}=-\frac{C'(\theta)}{g_\theta}<0. \nonumber
\end{align}
As $A$ decreases, $\theta^*$ (hence $\mu=\theta/(1-\theta)$) rises, the investment size falls and the
collateral rises, as claimed.

\answer{(balance sheet)} As $A$ decreases, the shadow cost $\mu$ rises, investment $I$ falls and
collateral $C$ rises: $dI/dA>0$, $dC/dA<0$.

\step{5}{Comparative statics in the agency cost $b$} Both $B$ and $p_H/\Delta p$ (with $p_H$ fixed)
enter only through $b=p_HB/\Delta p$, so it suffices to vary $b$. Now $b$ also enters $I(\theta,b)$; by
\eqref{t4e13-G}, $G_b=-I$, so
\begin{align}
  g_b &= G_b+G_I\frac{\partial I}{\partial b} \nonumber\\
      &= -I+\big(-(1-\theta)b\big)\big(-\theta Q\big) = -I+\theta(1-\theta)bQ, \label{t4e13-gb}
\end{align}
and by the implicit-function theorem $d\theta^*/db=-g_b/g_\theta$.
\emph{Investment.} Using \eqref{t4e13-Ith}, \eqref{t4e13-gth} and \eqref{t4e13-gb} (with $D\equiv G_CC'>0$),
\begin{align}
  \frac{dI}{db} &= \frac{\partial I}{\partial\theta}\frac{d\theta^*}{db}+\frac{\partial I}{\partial b}
               = (-bQ)\Big(-\frac{g_b}{g_\theta}\Big)+(-\theta Q)
               = -Q\Big(\theta-b\,\frac{g_b}{g_\theta}\Big) \nonumber\\
               &= -Q\,\frac{\theta g_\theta-bg_b}{g_\theta}
               = -Q\,\frac{\theta D+bI}{g_\theta}\ <0, \label{t4e13-Ib}
\end{align}
where the numerator is
$\theta g_\theta-bg_b=\theta(1-\theta)b^2Q+\theta D+bI-\theta(1-\theta)b^2Q=\theta D+bI>0$.
So $I$ falls with $B$ and with $p_H/\Delta p$, as claimed. \emph{Collateral.} $dC/db=C'(\theta)\,d\theta^*/db
=-C'(\theta)g_b/g_\theta$ has the sign of $-g_b$, and $-g_b>0$ iff $I>\theta(1-\theta)bQ$, i.e.\
(multiplying by $1/Q=p_H|R''|$) iff $p_H|R''|I>\theta(1-\theta)b$. By the definition of $\varepsilon$
and \eqref{t4e13-focI}, $p_H|R''|I=\varepsilon(I)\,p_HR'(I)=\varepsilon(I)(1+\theta b)$, so
\begin{align}
  \frac{dC}{db}>0 \iff \varepsilon(I)\,(1+\theta b)\ >\ \theta(1-\theta)\,b. \label{t4e13-cond}
\end{align}
This holds when $R$ is concave enough (for instance $\varepsilon\ge1$, since $\theta(1-\theta)b<1+\theta b$),
but it can fail when $R$ is close to linear. Counterexample: $R(I)=3I^{0.95}$ ($\varepsilon=0.05$),
$\phi(C)=\ln(1+C)$, $p_H=0.5$, $\Delta p=0.25$, $A=1$. For $B=0.5$ ($b=1$) the solution is
$\theta^*=0.3276$, $I=4.1203$, $C=1.9003$; for $B=2$ ($b=4$) it is $\theta^*=0.1238$, $I=0.3815$,
$C=0.3292$. Collateral \emph{falls} as the agency cost rises; at $b=1$, \eqref{t4e13-cond} reads
$0.05(1+0.3276)=0.0664>0.3276(1-0.3276)(1)=0.2203$, false. (The first best is not pledgeable in both
cases and $C\le BI/\Delta p$.)

\answer{(agency cost)} Investment always falls with the agency cost, $dI/db<0$. Collateral rises with
it iff $\varepsilon(I)(1+\theta b)>\theta(1-\theta)b$ (e.g.\ whenever $\varepsilon\ge1$); the statement's
claim for $C$ is not general, as the counterexample shows.

\intuition A weaker balance sheet tightens the financing constraint; the entrepreneur pays for funding
on both margins she controls, investing less (which saves outside funds) and pledging more collateral
(which adds pledgeable income at a deadweight cost). A higher agency cost also lowers investment, but
it has a second effect: each unit of investment now absorbs more pledgeable income, so the downsizing
itself releases funds. When output is nearly linear in $I$, downsizing is cheap and strong enough that
less collateral is needed. Only the balance-sheet comparative static is unambiguous.

\pitfalls
\begin{itemize}
  \item Dropping the indirect effect of collateral, the term $p_HC$: collateral also lets the cash
        reward $R_b$ fall by $C$ without breaking (IC$_b$).
  \item Treating $b$ as a pure shift of the constraint. It also enters the $I$-condition
        \eqref{t4e13-focI}, which is why the sign of $dC/db$ is ambiguous.
  \item Ignoring limited liability, $R_b=BI/\Delta p-C\ge0$: the substitution of collateral for
        reward is valid only while it holds.
\end{itemize}

\bookans Agrees with the book's answer (p.~591): the FOCs \eqref{t4e13-focI}--\eqref{t4e13-focC} and
``as the balance sheet deteriorates, $\mu$ increases, $I$ decreases, and $C$ increases''. The book's
answer does not claim that collateral rises with the agency cost (it says only that borrowing rises
when the agency cost falls). The statement's claim for $C$ holds only under \eqref{t4e13-cond}; the
claim for $I$ holds in general \eqref{t4e13-Ib}.

\refbib Tirole \S4.3.5, p.~170; \S3.4.2, pp.~127--128; answer p.~591.

% =============================================================================
\exsec{14. Dealing with asset substitution --- Tirole Exercise 7.3, pp.~323--324}

\asked (i) Without asset substitution: explain the threshold equation (1), give the
entrepreneur's expected utility, and describe the optimal contracts (at least at $A=\bar A$).
(ii) Now she can secretly run down maintenance (lower liquidation value $L$) to raise the success
probability. Explain condition (2), find the smallest $A$ for which this threat does no harm, and
describe the contract that achieves it.

\approach The fixed-investment model with a verifiable liquidation value (lecture Topic~4B, asset
substitution and remedies). Liquidation value is not subject to moral hazard, so all of it is
pledgeable; continuation income is pledgeable only up to $R-B/\Delta p$. A debt-like claim pays the
entrepreneur only in continuation, so she prefers to shift value from liquidation to continuation:
\emph{asset substitution}. To prevent it she must be paid something in liquidation, conditional on
good maintenance; that payment is a rent, which lowers \emph{pledgeable income} and raises the
minimum equity. We write (IR$_l$), (IC$_b$) and a no-substitution constraint, and minimise her rent.

\notation
\begin{itemize}
  \item $I>A$: fixed investment; $A$: cash; $x\in(0,1)$: probability of demand; with probability $1-x$
        there is no demand and the assets are liquidated for $L$.
  \item $R$: income in success; $p_H$, $p_L=p_H-\Delta p$, $B$: as in Section 3.2.
  \item $L_0$: liquidation value without asset substitution; $\tau(L)$: change in success probability
        when maintenance yields liquidation value $L$ ($\tau'<0$, $\tau''<0$, $\tau(L_0)=0$).
  \item $r_b$: entrepreneur's payment on liquidation; $R_b$: payment on continuation and success.
  \item $U_b^*$: entrepreneur's utility; $\bar A$: threshold (1); $A^*$: threshold with asset substitution.
\end{itemize}
\assume The statement's ``(hence, with probability $x$)'' for liquidation is read as probability $1-x$,
the probability of no demand. $L\ge0$ is chosen from $[0,\infty)$.

\step{1}{(i) Interpreting (1)} The liquidation value $L_0$ is verifiable and not subject to moral
hazard, so all of it can go to investors. On continuation only $R-R_b$ with $R_b\ge B/\Delta p$ can.
The left-hand side of (1) is the maximal expected pledgeable income,
\begin{align}
  \mathcal{P}_0=(1-x)L_0+xp_H\Big(R-\frac{B}{\Delta p}\Big), \label{t4e14-P0}
\end{align}
and (1) sets it equal to the investors' outlay at the threshold: financing is feasible iff
$A\ge\bar A=I-\mathcal{P}_0$. The side condition $L_0\le p_H(R-B/\Delta p)$ says that, when there is
demand, investors also prefer continuation (their maximal continuation claim beats $L_0$): there is
no conflict over the continuation decision.

\answer{(i), interpret (1)} The left side of (1) is the maximal expected pledgeable income (all of
$L_0$ on liquidation, $R-B/\Delta p$ on success after continuation); (1) defines $\bar A$ as the cash at
which it just covers the investors' outlay $I-A$.

\step{2}{(i) Utility and optimal contracts} Competitive investors break even, so the entrepreneur
receives the expected value created minus the investment, net of her own cash:
$U_b^*=\big[(1-x)L_0+xp_HR\big]-(I-A)-A$, i.e.
\begin{align}
  U_b^* = (1-x)L_0+xp_HR-I. \label{t4e14-U}
\end{align}
A contract $(r_b,R_b)$, liquidating iff there is no demand, is optimal iff
\begin{align}
  (1-x)(L_0-r_b)+xp_H(R-R_b) &= I-A, \label{t4e14-irl}\\
  R_b\ge\frac{B}{\Delta p},\quad r_b &\ge0. \label{t4e14-ic}
\end{align}
For $A>\bar A$ there is a continuum (the slack $A-\bar A$ can be paid as $r_b$ or as extra $R_b$). At
$A=\bar A$, subtracting (1), $I-\bar A=(1-x)L_0+xp_H(R-B/\Delta p)$, from \eqref{t4e14-irl},
\begin{align}
  -(1-x)\,r_b-xp_H\Big(R_b-\frac{B}{\Delta p}\Big)=0, \nonumber
\end{align}
and both terms are $\le0$ by \eqref{t4e14-ic}, so each is zero: $r_b=0$ and $R_b=B/\Delta p$. This is a debt
contract with face value $D=R-B/\Delta p\ge L_0$ (since $L_0\le p_H(R-B/\Delta p)\le D$). On liquidation
the firm defaults and investors take $L_0$; on success they are repaid $D$.

\answer{(i), utility and contracts} $U_b^*=(1-x)L_0+xp_HR-I$. Optimal contracts: any $(r_b,R_b)$ with
$r_b\ge0$, $R_b\ge B/\Delta p$ and binding \eqref{t4e14-irl}; at $A=\bar A$ the unique one is debt with
face $D=R-B/\Delta p$ ($r_b=0$, $R_b=B/\Delta p$).

\step{3}{(ii) Interpreting (2)} With maintenance level $L$ (and effort), the NPV is
\begin{align}
  \mathit{NPV}(L)=(1-x)L+x\big[p_H+\tau(L)\big]R-I. \label{t4e14-npv}
\end{align}
Differentiating, $\mathit{NPV}'(L)=(1-x)+x\tau'(L)R$ and $\mathit{NPV}''(L)=x\tau''(L)R<0$, so it is
strictly concave, and its first-order condition is $(1-x)+x\tau'(L)R=0$, i.e.\ (dividing by $x$)
$\tau'(L)R=-(1-x)/x$, which (2) says holds at $L_0$. So $L_0$ is the unique efficient maintenance level:
one more unit of liquidation value, worth $1-x$ in expectation, costs exactly $x|\tau'|R$ of expected
continuation income. $\tau(L_0)=0$ normalises the success probabilities so that $\mathit{NPV}(L_0)=U_b^*$.

\answer{(ii), interpret (2)} (2) says $L_0$ is the efficient maintenance level: it maximises the NPV,
and at $L_0$ the success probabilities are those of part (i).

\step{4}{(ii) The entrepreneur's maintenance choice} Under the step contract ($r_b(L)=r_b$ if $L\ge L_0$,
$0$ otherwise) and anticipating effort at stage 4 (her (IC$_b$), $(\Delta p)R_b\ge B$, does not depend
on $L$), she chooses $L$ to maximise
\begin{align}
  U(L)=(1-x)\,r_b(L)+x\big[p_H+\tau(L)\big]R_b. \label{t4e14-UL}
\end{align}
For $L<L_0$ she gets $0$ on liquidation and $\tau$ is decreasing, so the best deviation is $L=0$; any
$L>L_0$ lowers $\tau$ for no gain. So $L_0$ is chosen iff $U(L_0)\ge U(0)$:
\begin{align}
  (1-x)r_b+x\big[p_H+\tau(L_0)\big]R_b &\ge (1-x)(0)+x\big[p_H+\tau(0)\big]R_b \nonumber\\
  (1-x)r_b+xp_HR_b &\ge xp_HR_b+x\tau(0)R_b \qquad (\tau(L_0)=0) \nonumber\\
  (1-x)\,r_b &\ge x\,\tau(0)\,R_b. \tag{NS}\label{t4e14-ns}
\end{align}
A debt contract ($r_b=0$) violates \eqref{t4e14-ns}: the entrepreneur would run the assets down to gamble on
continuation.

\step{5}{(ii) Minimum $A$} Financing at the efficient $L_0$ requires the investors' breakeven
\eqref{t4e14-irl} with (IC$_b$) and \eqref{t4e14-ns}. The entrepreneur's expected rent
$(1-x)r_b+xp_HR_b$ is minimised by \eqref{t4e14-ns} at equality (any larger $r_b$ is pure rent);
substituting $(1-x)r_b=x\tau(0)R_b$, the rent is $x\tau(0)R_b+xp_HR_b=x\big[p_H+\tau(0)\big]R_b$,
increasing in $R_b$, so $R_b=B/\Delta p$, the lowest value (IC$_b$) allows. Then
\begin{align}
  r_b=\frac{x\,\tau(0)}{1-x}\,\frac{B}{\Delta p}. \label{t4e14-rb}
\end{align}
Substituting $R_b=B/\Delta p$ and \eqref{t4e14-rb} into \eqref{t4e14-irl} at equality,
\begin{align}
  I-A^* &= (1-x)\big(L_0-r_b\big)+xp_H\Big(R-\frac{B}{\Delta p}\Big) \nonumber\\
  A^* &= \Big[I-(1-x)L_0-xp_H\Big(R-\frac{B}{\Delta p}\Big)\Big]+(1-x)\,r_b \nonumber\\
      &= \bar A+(1-x)\,r_b=\bar A+x\,\tau(0)\,\frac{B}{\Delta p}, \label{t4e14-Astar}
\end{align}
using (1) for the bracket and \eqref{t4e14-rb} for $(1-x)r_b$.
For $A\ge A^*$, asset substitution is innocuous: the firm is financed, maintains $L_0$ and the
entrepreneur gets $U_b^*$. For $A\in[\bar A,A^*)$ the efficient contract is not fundable. Illustration:
$x=0.6$, $p_H=0.8$, $p_L=0.4$, $B=1$, $R=10$, $L_0=3$, $I=5.5$, $\tau(L)=(9-L^2)/90$ (which satisfies
(2): $\tau'(3)R=-(6/90)(10)=-2/3=-(0.4/0.6)$): $U_b^*=0.4(3)+0.6(0.8)(10)-5.5=1.2+4.8-5.5=0.5$,
$B/\Delta p=2.5$, $\bar A=5.5-1.2-0.6(0.8)(10-2.5)=5.5-1.2-3.6=0.7$, $\tau(0)=9/90=0.1$,
$r_b=(0.6)(0.1)(2.5)/(0.4)=0.375$, $A^*=0.7+(0.4)(0.375)=0.85$.

\step{6}{(ii) The optimal contract} It is debt plus a reward for good maintenance: the entrepreneur
keeps $B/\Delta p$ on success and receives $r_b$ on liquidation, but only if the assets are in good
condition (a severance payment conditional on $L\ge L_0$). Giving her a stake in the liquidation
value counters her bias toward the continuation state, where her debt-like claim pays.

\answer{(ii), minimum $A$ and contract} $A^*=\bar A+x\,\tau(0)\,B/\Delta p$. The contract is debt
($R_b=B/\Delta p$ on success) plus a payment $r_b=x\tau(0)B/[(1-x)\Delta p]$ on liquidation, made only if
the assets were maintained at $L_0$.

\intuition A claim that pays only in success makes the entrepreneur care only about continuation
profit, so she cuts maintenance to raise the success probability: risk shifting at the expense of
investors, who own the liquidation value. To stop it, she must share in the liquidation value. That
share is an extra rent, so it lowers pledgeable income and raises the equity needed, by the expected
gain from the most aggressive substitution, $x\,\tau(0)B/\Delta p$.

\pitfalls
\begin{itemize}
  \item Reading ``liquidation \ldots\ (hence, with probability $x$)'' literally; liquidation happens when
        there is no demand, with probability $1-x$ (Figure~7.16).
  \item Keeping the debt contract of part (i) in part (ii). With $r_b=0$ the no-substitution constraint
        fails and she runs the assets down to $L=0$.
  \item Paying $r_b$ on liquidation unconditionally. A payment that does not depend on $L$ leaves her
        maintenance incentive unchanged; it must be conditional on $L\ge L_0$.
\end{itemize}

\bookans Matches the book's answer (p.~602): the same interpretation of (1), $U_b^*$, the debt contract
at $A=\bar A$, the step-function contract with $R_b=B/\Delta p$, the no-substitution inequality
$(1-x)r_b+xp_HB/\Delta p\ge x[p_H+\tau(0)]B/\Delta p$ and the threshold
$I-A^*=(1-x)(L_0-r_b)+xp_H(R-B/\Delta p)$. The statement's ``with probability $x$'' for liquidation
looks like a typo for $1-x$.

\refbib Tirole \S3.2.2, pp.~116--117; \S7.2.3 ``Effort and Risk Taking'', p.~314; answer p.~602.

% =============================================================================
\exsec{15. A borrower-friendly bankruptcy court --- Tirole Exercise 4.9, pp.~190--191}

\asked (i) Argue that investors should get the short-term profit and the liquidation value, that
the entrepreneur's continuation reward should be exactly $B/\Delta p$, and explain assumptions (1)
and (2). (ii) Write the borrower's optimisation problem and the equation for the continuation cutoff
$r^*$. (iii) Show that short-term debt $d=r^*$ implements the contract (default $=$ liquidation) and
how $d$ depends on her cash $A$. (iv) A court splits the liquidation proceeds equally with the
borrower; show that if $r^*>\hat r$ she loses access to financing. (v) If $r^*<\hat r$, show that the
court either blocks financing or makes bankruptcy more likely, hurting the borrower but not lenders.

\approach Lecture Topic~4B, bankruptcy design. The fixed-investment model with an intermediate
verifiable profit $r$ and a liquidation option. Everything not subject to moral hazard ($r$, $L$) is
\emph{pledgeable income} at no incentive cost; continuation income is pledgeable only up to
$\rho_0(r)$. When the project cannot be financed with continuation in every state, the borrower buys
funding by accepting inefficient liquidation where it raises pledgeable income most cheaply, which
gives a cutoff rule. We write the Lagrangian, maximise it state by state, and use the investors'
breakeven condition to find $r^*$; the court changes investors' payoff from liquidation from $L$ to
$L/2$.

\notation
\begin{itemize}
  \item $r\in[0,\bar r]$: verifiable short-term profit, density $f>0$, c.d.f.\ $F$; $L>0$: liquidation value.
  \item $p_H(r)$, $p_L(r)=p_H(r)-\Delta p$: success probabilities ($p_H'>0$); $R$: final profit in success;
        $B$: private benefit.
  \item $\rho_1(r)=p_H(r)R$: continuation value; $\rho_0(r)=p_H(r)(R-B/\Delta p)$: pledgeable
        continuation value; both increasing in $r$.
  \item $\ell(r)\in\{0,1\}$: liquidation policy; $R_b(r)$: reward on continuation and success;
        $r^*$: optimal cutoff; $d$: short-term debt.
  \item $\hat r$: $\rho_0(\hat r)=L/2$; $r^{**}$: cutoff under the borrower-friendly court.
  \item $P(c)$, $P_c(c)$: investors' expected income with cutoff $c$ under an enforcing / borrower-friendly court.
\end{itemize}
\assume The court's equal split of $L$ cannot be undone by contract.

\step{1}{(i) What goes to investors; $R_b=B/\Delta p$; conditions (1)--(2)} $r$ and $L$ are not subject to
moral hazard, so giving them to investors raises pledgeable income without hurting incentives. On
continuation, (IC$_b$) is $[p_H(r)-p_L(r)]R_b\ge B$, i.e.\ $R_b\ge B/\Delta p$; any excess is wasted
pledgeable income, which would force more liquidation (proved in Step 2). Condition (1), $\rho_1(r)>L$,
says continuation is always efficient. Condition (2): always liquidating raises enough,
$E[r]+L>I-A$, but always continuing does not, $E[r+\rho_0(r)]<I-A$. Financing therefore needs
inefficient liquidation in some states.

\answer{(i)} $r$ and $L$ go to investors because pledging them costs no incentives; $R_b=B/\Delta p$
because any larger reward wastes pledgeable income. (1): continuation is always efficient. (2):
liquidating always would raise enough funds, continuing always would not, so some inefficient
liquidation is the price of financing.

\step{2}{(ii) Programme and the cutoff} With binding breakeven the borrower's utility is the NPV:
\begin{align}
  \max_{\ell(\cdot),R_b(\cdot)}\; & E[r]+\int_0^{\bar r}\big[\ell(r)L+(1-\ell(r))\rho_1(r)\big]f(r)\,dr-I \label{t4e15-obj}\\
  \text{s.t. }\; & E[r]+\int_0^{\bar r}\big[\ell(r)L+(1-\ell(r))p_H(r)(R-R_b(r))\big]f(r)\,dr\ \ge\ I-A, \label{t4e15-irl}\\
  & R_b(r)\ge B/\Delta p. \label{t4e15-ic}
\end{align}
Let $\mu$ be the multiplier of \eqref{t4e15-irl}. $\mu>0$: the unconstrained optimum, $\ell\equiv0$ by (1),
violates \eqref{t4e15-irl} by (2). $R_b$ does not enter \eqref{t4e15-obj} and lowers \eqref{t4e15-irl}, so
$R_b(r)=B/\Delta p$ and the investors' continuation income is $p_H(r)(R-B/\Delta p)=\rho_0(r)$. The
Lagrangian is then
\begin{align}
  \mathcal L &= (1+\mu)E[r]-I-\mu(I-A)+\int_0^{\bar r}\Big\{\ell(r)\,(1+\mu)L
               +\big(1-\ell(r)\big)\big[\rho_1(r)+\mu\rho_0(r)\big]\Big\}f(r)\,dr, \nonumber
\end{align}
which is linear in $\ell(r)$ state by state. Maximising the integrand pointwise: liquidate iff
$(1+\mu)L>\rho_1(r)+\mu\rho_0(r)$, i.e.
\begin{align}
  \mu\big[L-\rho_0(r)\big]\ >\ \rho_1(r)-L. \label{t4e15-point}
\end{align}
The left-hand side decreases and the right-hand side increases in $r$, so the policy is a cutoff:
liquidate iff $r<r^*$. At the cutoff \eqref{t4e15-point} holds with equality and $\rho_1(r^*)>L$, so
$\mu[L-\rho_0(r^*)]=\rho_1(r^*)-L>0$ and, with $\mu>0$, $\rho_0(r^*)<L$.

\step{3}{(ii) The equation for $r^*$} With cutoff $c$, investors receive
\begin{align}
  P(c)=E[r]+L\,F(c)+\int_c^{\bar r}\rho_0(r)f(r)\,dr, \label{t4e15-P}
\end{align}
and the NPV, $E[r]+LF(c)+\int_c^{\bar r}\rho_1f\,dr-I$, falls with $c$ by (1). So $r^*$ is the smallest
root of
\begin{align}
  E[r]+L\,F(r^*)+\int_{r^*}^{\bar r}\rho_0(r)f(r)\,dr = I-A. \label{t4e15-rstar}
\end{align}
By (2), $P(0)=E[r+\rho_0(r)]<I-A<E[r]+L=P(\bar r)$, so by continuity $r^*\in(0,\bar r)$ (the
statement's ``$(0,r)$'' should read $(0,\bar r)$).

\answer{(ii)} $\max_{\ell(\cdot),R_b(\cdot)}$ \eqref{t4e15-obj} s.t.\ \eqref{t4e15-irl}--\eqref{t4e15-ic};
continue iff $r\ge r^*$, with $r^*$ the smallest root of
$E[r]+LF(r^*)+\int_{r^*}^{\bar r}\rho_0(r)f(r)\,dr=I-A$.

\step{4}{(iii) Short-term debt and the effect of $A$} Implementation: short-term debt $d=r^*$ due at
date 1, any excess $r-d$ also paid to investors (for example as a payout to outside shareholders), and a
long-term claim of $R-B/\Delta p$ on success. If $r\ge d$ the debt is paid and the firm continues; if
$r<d$ it defaults, creditors take control and liquidate (``bankruptcy''). They do not renegotiate:
for $r<r^*$, $\rho_0(r)<\rho_0(r^*)<L$, and the borrower has no cash to bribe them. By Leibniz's
rule applied to \eqref{t4e15-P},
\begin{align}
  \frac{dP}{dc}=L\,f(c)-\rho_0(c)f(c)=\big[L-\rho_0(c)\big]f(c). \nonumber
\end{align}
Differentiating \eqref{t4e15-rstar}, $P(r^*)=I-A$, with respect to $A$:
\begin{align}
  \big[L-\rho_0(r^*)\big]f(r^*)\,dr^* &= -dA \nonumber\\
  \frac{dr^*}{dA} &= -\frac{1}{[L-\rho_0(r^*)]f(r^*)}\ <0. \label{t4e15-drA}
\end{align}
Short-term debt falls with initial equity: a weaker balance sheet needs more pledgeable income, so the
firm must accept bankruptcy in more states.

\answer{(iii)} Short-term debt $d=r^*$ (default and liquidation iff $r<d$), plus a long-term claim
$R-B/\Delta p$ on success, implements the optimum; $dd/dA=dr^*/dA=-1/\{[L-\rho_0(r^*)]f(r^*)\}<0$.

\step{5}{(iv) Borrower-friendly court with $r^*>\hat r$} Liquidation now yields investors only $L/2$. For
any policy, investors' income in state $r$ (besides $r$) is at most $\max\{L/2,\rho_0(r)\}$. Compare
with the enforcing-court contract state by state. For $r<r^*$, $\max\{L/2,\rho_0(r)\}<L$, since
$L/2<L$ and $\rho_0(r)<\rho_0(r^*)<L$. For $r\ge r^*>\hat r$, $\rho_0(r)>\rho_0(\hat r)=L/2$, so the
maximum is $\rho_0(r)$, the same as before. Hence
\begin{align}
  E[r]+\int_0^{\bar r}\max\Big\{\frac L2,\rho_0(r)\Big\}f(r)\,dr\ <\ E[r]+LF(r^*)+\int_{r^*}^{\bar r}\rho_0f\,dr = I-A, \label{t4e15-iv}
\end{align}
using \eqref{t4e15-rstar}. No policy lets investors break even: the borrower loses access to financing.

\answer{(iv)} With $r^*>\hat r$, the investors' maximal expected income under the court is below
$I-A$: no contract is financed.

\step{6}{(v) Court with $r^*<\hat r$} Cutoff policies are again optimal (same argument as
\eqref{t4e15-point} with $L/2$ for investors). Investors' income with cutoff $c$ is
\begin{align}
  P_c(c)=E[r]+\frac L2F(c)+\int_c^{\bar r}\rho_0(r)f(r)\,dr, \qquad P_c'(c)=\Big[\frac L2-\rho_0(c)\Big]f(c), \label{t4e15-Pc}
\end{align}
which rises up to $\hat r$ and falls after it, so its maximum is $P_c(\hat r)$. If $P_c(\hat r)<I-A$, there
is no financing. Otherwise the borrower picks the smallest root $r^{**}$ of $P_c(c)=I-A$. For
$0<c\le r^*$, $P_c(c)<P(c)\le I-A$ by \eqref{t4e15-P} and \eqref{t4e15-rstar}, so $r^{**}\in(r^*,\hat r]$:
bankruptcy is more likely. Lenders break even in either regime, so they are not hurt. The borrower gets
the NPV, which falls with the cutoff (Step 3), so $\mathit{NPV}(r^{**})<\mathit{NPV}(r^*)$: she is hurt.
More directly, with an enforcing court she could have written the $L/2$ split into her contract; the
friendly court only removes options.

\answer{(v)} If $P_c(\hat r)<I-A$ there is no financing; otherwise the cutoff rises to
$r^{**}\in(r^*,\hat r]$ (more bankruptcy). Lenders still break even; the borrower's utility, the NPV,
falls.

\step{7}{Numerical illustration} $r\sim U[0,1]$, $p_H(r)=0.5+0.4r$, $\Delta p=0.2$, $R=10$, $B=1.2$
($R-B/\Delta p=4$), $L=4.5$. Then $\rho_1(r)=5+4r>L$, $\rho_0(r)=2+1.6r$, $E[r]=0.5$,
$P(0)=E[r+\rho_0]=3.3$, $P(1)=E[r]+L=5$, so (2) holds for $I-A\in(3.3,5)$; $\hat r=(2.25-2)/1.6=0.15625$,
and the court's maximal income is $P_c(\hat r)=3.319531$. With $f=1$ on $[0,1]$, \eqref{t4e15-P} and
\eqref{t4e15-Pc} integrate to
\begin{align}
  P(c) &= 0.5+4.5c+2(1-c)+0.8(1-c^2)=3.3+2.5c-0.8c^2, \nonumber\\
  P_c(c) &= 0.5+2.25c+2(1-c)+0.8(1-c^2)=3.3+0.25c-0.8c^2, \nonumber
\end{align}
so $P_c(\hat r)=3.3+0.25(0.15625)-0.8(0.15625)^2=3.3+0.0390625-0.0195313=3.319531$.
(iv) $I-A=4$: $P(r^*)=4$ gives $0.8r^{*2}-2.5r^*+0.7=0$, smallest root
$r^*=(2.5-\sqrt{6.25-2.24})/1.6=0.3109>\hat r$; $\rho_0(r^*)=2+1.6(0.3109)=2.4975<L$; and $3.319531<4$:
no financing. (v) $I-A=3.31$: $2.5r^*-0.8r^{*2}=0.01$ gives $r^*=0.0040<\hat r$; under the court
$0.25r^{**}-0.8r^{**2}=0.01$ gives $r^{**}=(0.25-\sqrt{0.0625-0.032})/1.6=0.0471$, and the NPV loss is
$\int_{r^*}^{r^{**}}(\rho_1-L)\,dr=\int_{r^*}^{r^{**}}(0.5+4r)\,dr=0.0260$.

\intuition Liquidation is the investors' collateral: it is what they get when the firm cannot pay. A court
that hands half of it to the borrower lowers the value of default to creditors. To keep them whole the
contract must liquidate more often, which destroys value, or financing breaks down. Lenders still break
even, so the whole cost of the borrower-friendly rule falls on the borrower it was meant to protect.

\pitfalls
\begin{itemize}
  \item Liquidating whenever $\rho_0(r)<L$. Liquidation is always inefficient here (1); it is used only
        as far as needed to break even, so the cutoff comes from the binding breakeven, not from
        $\rho_0=L$.
  \item Concluding that a court that hands the borrower half of $L$ helps her. Lenders must still
        break even, so the borrower pays for the transfer through more liquidation or no financing.
  \item Signing $dr^*/dA$ wrong: more cash means less pledgeable income is needed, so fewer
        liquidations and \emph{less} short-term debt.
\end{itemize}

\bookans Matches the book's answer (pp.~586--587): the reasons for pledging $r$ and $L$, $R_b=B/\Delta p$,
the programme and $r^*$ as the lowest root of breakeven with $L\ge\rho_0(r^*)$, the debt implementation
$d=r^*$, $dr^*/dA<0$, and (iv)--(v) as read from the book's Figures 2--3 (bankruptcy on $[0,r^{**}]$ with
$r^*<r^{**}\le\hat r$, lenders not hurt, borrower hurt). The statement's ``$(0,r)$'' is a typo for $(0,\bar r)$.

\refbib Tirole \S3.2.2, pp.~116--117; Exercise 4.9, pp.~190--191; answer pp.~586--587.

\end{document}
